\documentclass[11pt,a4paper]{article}
\usepackage[margin=23mm]{geometry}
\usepackage[T1]{fontenc}
\usepackage[utf8]{inputenc}
\usepackage{lmodern}
\usepackage{amssymb,amsmath,amsthm,mathtools}
\usepackage[unicode,colorlinks=true,linkcolor=blue,urlcolor=blue]{hyperref}
\usepackage{microtype}
\setlength{\emergencystretch}{2em}
\newcommand{\R}{\mathbb R}
\newcommand{\N}{\mathbb N}
\newcommand{\eps}{\varepsilon}
\newcommand{\carrier}[1]{\left|#1\right|}
\newcommand{\typeOf}{\operatorname{type}}
\newtheorem{definition}{Definition}[section]
\newtheorem{lemma}[definition]{Lemma}
\newtheorem{theorem}[definition]{Theorem}
\newtheorem{proposition}[definition]{Proposition}
\newtheorem{corollary}[definition]{Corollary}
\theoremstyle{remark}
\newtheorem{remark}[definition]{Remark}

\newcommand{\lean}[1]{}
\newcommand{\uses}[1]{}
\newcommand{\proves}[1]{}
\newcommand{\discussion}[1]{}
\newcommand{\mathlibok}{}
\newcommand{\leanok}{\textit{Lean kernel checked.}\ }
\newcommand{\notready}{\textbf{Open construction dependency.}\ }

\title{A Lean Blueprint for Substitution and Hierarchical Tilings}
\author{Formalization design and proof repair}
\date{12 October 2026, review revision}
\begin{document}
\maketitle
\tableofcontents
\section{Scope, status, and exact models}

This blueprint covers all five conclusions of Theorems 1 and 2 of
\href{https://arxiv.org/abs/2606.25005v2}{arXiv:2606.25005v2}.
It distinguishes three statuses: an ordinary mathematical proof supplied
below, a Lean declaration checked in the small header kernel, and a
geometric construction obligation still requiring a complete proof.
Only nodes corresponding exactly to the checked header kernel carry
\texttt{leanok}. All other Lean names are planned names.
In particular, the universal main theorems are goals until the outstanding
geometric instantiation is established. A certificate theorem is an
intermediate theorem, not a replacement for constructing its certificate.

The new construction changes the boundary control colour in the hierarchical
case. The crown-cover proof uses new positive constants. These changes do
not add convexity, primitivity, recognizability, or lattice assumptions.
We interpret polygonal tiles as compact simple polygonal disks. Collinear
listed vertices are permitted. Rotations are represented by distinct forms
before the construction; only translations act on placed tiles.

\begin{definition}[Polygonal disk]\label{def:polygon}
\lean{Tilings.Geometry.PolygonalDisk}
A polygonal disk is a compact regular-closed set $B\subset\R^2$ with
nonempty interior and a finite simple polygonal boundary. It carries
indexed boundary vertices and nonzero segments. Consecutive collinear
vertices may be retained as data. Its interior is path connected.
\end{definition}

For Lean, use $\operatorname{EuclideanSpace}\R(\operatorname{Fin}2)$.
The usual product metric on $\R\times\R$ is not silently substituted for
the Euclidean metric in the distance and sine estimates. A vertex-only
\texttt{Polygon} interface is not yet a filled polygonal disk; the bridge
between boundary data and the carrier is a separate development.

\begin{definition}[Finite protoset and placements]\label{def:protoset}
\lean{Tilings.Protoset}
\uses{def:polygon}
A protoset is a finite type $P$ with a polygonal disk $B_p$ for each $p\in P$.
A placed tile is $(p,v)$, with carrier $v+B_p$. Colours belong to the form,
so congruent coloured forms remain distinct.
\end{definition}

\begin{definition}[Placed tiles and patches]\label{def:placed}
\lean{Tilings.PlacedTile,Tilings.Patch}
\uses{def:protoset}
Write $|A|$ for the carrier of a placed tile and $|F|$ for the union of the
carriers in a patch. A patch is a finite set of placed tiles with pairwise
disjoint interiors. Translations act on the anchor coordinates.
\end{definition}

\begin{definition}[Plane tilings and side-to-side]\label{def:tiling}
\lean{Tilings.PlaneTiling,Tilings.SideToSide}
\uses{def:placed}
A plane tiling $T$ is a set of placed tiles whose carriers cover $\R^2$ and
whose interiors are pairwise disjoint. It is side-to-side if two tile sides
with positive-length overlap coincide as whole segments. Local finiteness
is a derived property, not an extra hypothesis of the final theorem.

\emph{Plane coverage is part of every theorem statement below.} The
paper's introductory word ``tiling'' (main.tex, line 92) also admits an
arbitrary nonoverlapping set; that usage is kept here only for finite
patches. Properness of a mere nonoverlapping set is not the notion used:
without coverage a single inert decorated tile would be a proper tiling
projecting outside the family, the surjectivity device of
Lemma~\ref{lem:glues} and the reverse direction of the collar conversion
would fail, and nontrivial sides could lack an opposite neighbour.
\end{definition}

\begin{lemma}[Local finiteness and countability]\label{lem:localfinite}
\lean{Tilings.PlaneTiling.locallyFinite,Tilings.PlaneTiling.countable}
\uses{def:tiling}
Every plane tiling by a finite protoset is locally finite and countable.
\end{lemma}
\begin{proof}
Choose an open disk in the interior of each form. The minimum of these
finitely many positive radii is $r>0$, and the maximum of the form diameters
is $D<\infty$. If a tile meets a compact set $K$, its chosen interior disk
centre lies in the closed $D$-neighbourhood of $K$. Centres belonging to
distinct tiles are at distance at least $2r$, since their open interior
disks are disjoint. A compact metric set has only finitely many
$2r$-separated points, by a finite cover with disks of radius $r/2$.
This proves local finiteness.
For countability choose, for each tile, a rational point in its nonempty
open interior. Chosen rational points for different tiles are distinct.
The injection into $\mathbb Q^2$ proves countability. This proof does not
require rational polygon coordinates.
\end{proof}

\begin{definition}[Substitution]\label{def:substitution}
\lean{Tilings.Substitution}
\uses{def:placed}
A substitution $\tau$ has a common scale $\theta>1$ and a finite patch
$\tau p$ with support $\theta B_p$ for each $p\in P$. On a placed tile
$(p,v)$ it produces the translates by $\theta v$ of $\tau p$.
Iterates are defined recursively. A supertile is a translate of $\tau^n p$,
$n\in\N$. A composition is an inverse decomposition of a plane tiling.
Decorated substitutions are allowed to be relations rather than functions.
\end{definition}

\begin{lemma}[Supports and finite parents]\label{lem:finiteparents}
\lean{Tilings.Substitution.iterate_support,Tilings.Substitution.finite_parents}
\uses{def:substitution}
The support of $\tau^n p$ is $\theta^n B_p$. A fixed placed child has only
finitely many possible placed parents under $\tau$.
\end{lemma}
\begin{proof}
The support identity follows by induction, using the exact support identity
for every child and distribution of a homothety over a union.
A possible parent is determined by its form and the child position in that
form's finite substitution patch. If the child's anchor is $v$ and the
chosen child offset is $w$, the parent's anchor is $(v-w)/\theta$.
There are finitely many such pairs of form and position.
\end{proof}

\begin{definition}[Substitution and hierarchical families]\label{def:hierarchy}
\lean{Tilings.IsSubstitutionTiling,Tilings.IsHierarchical,Tilings.IsSideToSideHierarchical}
\uses{def:tiling,def:substitution}
A plane tiling is a substitution tiling if every finite patch in it occurs,
up to translation, in a supertile. It is hierarchical if there are plane
tilings $T_0=T,T_1,\ldots$ with $\tau T_{n+1}=T_n$.
It is side-to-side hierarchical if every $T_n$ is side-to-side.
\end{definition}

\begin{lemma}[Substitution tilings admit substitution compositions]
\label{lem:substitution-hierarchy}
\lean{Tilings.SubstitutionTiling.exists_composition}
\uses{def:hierarchy,lem:finiteparents,lem:localfinite}
Every substitution plane tiling has a composition which is again a
substitution plane tiling.
\end{lemma}
\begin{proof}
Enumerate its tiles $A_1,A_2,\ldots$. Each initial finite set is contained in
a supertile of positive order: enlarge the patch by sufficiently many
tiles if necessary to exclude order zero. Choose the previous-level
supertile as a parent cover. By the finite-parent lemma, thin this sequence
to stabilize the parent of $A_1$, then of $A_2$, and so forth. The diagonal
sequence yields parents $B_i$ such that every finite set of the $B_i$ is
contained in a supertile. Thus they have disjoint interiors and form a
substitution patch at every finite scale. Their decompositions contain
every $A_i$. No extra child can occur: a child has nonempty interior, and
the plane is already covered by $T$ with disjoint interiors. Consequently
the decomposition is exactly $T$; scaling gives coverage of the parent
plane. Iterate this existence statement using dependent choice.
The Lean proof may equivalently use a finite inverse system of partial
parent assignments and K\"onig's lemma.
\end{proof}

\begin{lemma}[Powers and hierarchical tilings]\label{lem:hierarchy-step}
\lean{Tilings.hierarchical_power_iff}
\uses{def:hierarchy}
$\tau$-hierarchy and $\tau^m$-hierarchy coincide for $m>0$.
Side-to-side hierarchy requires a separate check on the inserted levels.
\end{lemma}
\begin{proof}
Take every $m$th level in one direction. In the other, between two levels
insert the $m-1$ intermediate decompositions of the upper tiling. Exact
supports and disjoint interiors ensure that these are plane tilings.
Side-to-side at the two endpoints does not by itself imply side-to-side
at the intermediate levels, which is why the alignment tuples are needed.
\end{proof}

\begin{definition}[Crowns, FLC, and STS]\label{def:crown}
\lean{Tilings.Crown,Tilings.AllowedCrown}
\uses{def:placed,def:substitution,lem:localfinite}
The full pointed crown at a listed tile vertex $v$ consists of all tiles
whose closed carriers contain $v$, with distinguished centre $v$.
A crown is allowed if it occurs in a supertile. Equality up to translation
preserves the distinguished centre.
\end{definition}

\begin{definition}[Finite local complexity]\label{def:flc}
\lean{Tilings.Substitution.FLC}
\uses{def:crown}
FLC means that the set of allowed full pointed crowns, including incomplete
crowns on supertile boundaries, is finite up to translation.
\end{definition}

\begin{definition}[STS hypothesis]\label{def:sts}
\lean{Tilings.Substitution.STS}
\uses{def:tiling,def:substitution}
STS means that every supertile is side-to-side.
\end{definition}


\begin{definition}[Hereditary listed sides]\label{def:hereditary}
\lean{Tilings.Substitution.HereditarySides}
\uses{def:substitution,def:sts}
Every image of a listed prototile vertex is a listed child vertex, and each
scaled parent listed side is a union of complete child listed sides.
This is an extra core hypothesis, not a consequence of STS: an irrational
dummy cut on unit-square sides is not hereditary for a $2\times2$ subdivision.
Finite refinement retaining such a cut cannot repair heredity, since it
would contain the infinite orbit $\{2^n a\}\bmod1$.
The final main theorems remove this extra hypothesis explicitly.
\end{definition}

\begin{lemma}[STS gives FLC]\label{lem:sts-flc}
\lean{Tilings.Substitution.STS.finite_crowns}
\uses{def:sts,def:flc}
Under STS the allowed pointed crowns form a finite set up to translation.
\end{lemma}
\begin{proof}
At a vertex, a tile in the crown is fixed by its form and the listed vertex
identified with the centre. In a face-to-face polygonal patch there is no
vertex strictly inside another tile side with an incident open sector:
the sector on that side would force a partial side overlap. The number of
incident tiles is bounded by the smallest positive angle among the finite
forms. Hence a crown is a subset of a finite list of positioned forms.
Incomplete crowns are subsets of the same list.
The proof of the no-partial-side assertion uses that a supertile fills its
simple polygonal carrier, rather than an arbitrary disconnected patch.
\end{proof}

\begin{definition}[Bounded local rules and finite decorations]\label{def:localrule}
\lean{Tilings.LocalRule,Tilings.Decoration,Tilings.Proper}
\uses{def:tiling}
A local rule is a possibly infinite set $\mathcal R$ of finite forbidden
patches, together with $D<\infty$ bounding all their support diameters.
A full plane tiling is proper if it contains no translated occurrence of
a member of $\mathcal R$. A decoration is a finite decorated form type,
with a surjection to the original forms preserving carriers after anchor
normalization. Its induced projection acts tilewise.
\end{definition}

\begin{definition}[Soficity]\label{def:sofic}
\lean{Tilings.IsSofic,Tilings.IsSideToSideSofic}
\uses{def:localrule}
A family is sofic if it is exactly the projections of the proper plane
tilings for a finite decoration and a bounded local rule.
It is side-to-side sofic if the rule is equality of colours on complete
matching sides, together with side-to-side adjacency.
\end{definition}

This is the continuous geometric notion in the paper. It is not replaced
by a $\mathbb Z^2$ subshift of finite type. The forbidden set may be infinite;
the decoration alphabet and the radius are finite. Both quantifiers in the
definition of soficity range over plane tilings only, as fixed in
Definition~\ref{def:tiling}.

\begin{lemma}[Finite endofunction has an idempotent power]
\label{lem:finite-idempotent}
\lean{Tilings.FiniteMap.exists_idempotent_iterate}
For a function $c$ on a finite set there is $\ell>0$ with $c^{2\ell}=c^\ell$.
\end{lemma}
\begin{proof}
Finitely many functions imply $c^a=c^{a+b}$ for some $a\geq0,b>0$.
Then $c^{a+i}=c^{a+b+i}$ for every $i\geq0$. Choose a positive multiple
$\ell$ of $b$ with $\ell\geq a$. Both $\ell$ and $2\ell$ lie in the periodic
tail and have the same residue modulo $b$, giving the equality.
Restricting the image of $c^\ell$ makes every retained image point fixed
by $c^\ell$. This is the central-form stabilization used below.
\end{proof}

\section{A corrected uniform crown-cover lemma}
\label{sec:corrected-crowns}

This section gives an ordinary mathematical proof of the uniform
crown-cover statement used in the FLC reduction.  All Lean declarations
attached to the nodes below are planned declarations.  None of these nodes
is marked as formally proved.  The proof uses full pointed crowns, including
incomplete crowns on supertile boundaries, and replaces both problematic
estimates in the Appendix: the endpoint estimate by a minimal prototile side
length, and the sine of an exterior angle exceeding $\pi$.

Throughout, polygons are compact simple polygonal disks with nonempty
interior, positive-length edges, and carrier equal to the closure of their
interior.  They may be nonconvex.  All distances and diameters are Euclidean.
Write $B(v,r)$ for the open ball of radius $r$ about $v$.

\begin{definition}[Pointed crown catalogue and sectors]
\label{def:crown-sector}
\lean{Tilings.Crowns.pointedCatalogue}
\uses{def:polygon,def:protoset,def:placed,def:tiling,def:crown,def:flc}
A \emph{full pointed crown} at a tile vertex $v$ consists of $v$ together
with every tile whose closed carrier contains $v$.  Fix a finite catalogue
$\mathcal C$ of such pointed configurations, up to translation.  Every crown
to which the following results are applied must occur in this catalogue.
For supertiles the catalogue includes incomplete boundary crowns, rather
than only complete crowns from plane tilings.

If a polygon $A$ contains $v$, let $T_vA$ be its closed local tangent cone:
the whole plane when $v$ is an interior point, its inward halfplane when
$v$ lies strictly inside a straight boundary edge, and its interior wedge
when $v$ is a genuine corner.  Collinear listed vertices use the halfplane
case.  For a pointed crown $K$ set
\[
  \Omega(K)=\bigcup_{A\in K}T_vA,
  \qquad \operatorname{Sec}(K,r)=B(v,r)\cap\Omega(K).
\]
The union here includes tiles for which the centre $v$ lies on an edge
rather than at one of their own listed vertices.
\end{definition}

\begin{lemma}[Local polygon cones]
\label{lem:polygon-local-cone}
\lean{Tilings.Crowns.polygonLocalCone}
\uses{def:polygon,def:placed,def:crown-sector}
For every polygon $A$ and every $v\in A$, there is a number $r>0$ such that
\[
  A\cap B(v,r)=T_vA\cap B(v,r).
\]
\end{lemma}
\begin{proof}
For an interior point take a ball contained in the interior.  At a boundary
point, merge consecutive collinear boundary segments for this local
argument.  The compact union of boundary edges not incident to the point
has positive distance from it.  Decrease the radius further to avoid the
other endpoints of the one or two incident edges.  The boundary inside the
ball is consequently one straight line segment, or the two rays at a
genuine corner.  The simple polygonal disk lies locally on its prescribed
interior side, giving respectively the inward halfplane or interior wedge.
Keeping collinear vertices in the original vertex list changes neither
the disk nor this local description.
\end{proof}

\begin{lemma}[A covered tile belongs to the covering patch]
\label{lem:covered-tile-belongs}
\lean{Tilings.Patch.mem_of_carrier_subset_union}
\uses{def:polygon,def:placed}
Let $F$ be a finite patch and let $X$ be a tile whose interior is disjoint
from the interior of every tile of $F$ other than $X$ itself. If the
carrier of $X$ is contained in $\bigcup F$, then $X\in F$.
\end{lemma}
\begin{proof}
If $X\notin F$, the nonempty open set $\operatorname{int}X$ is covered by
the finitely many carriers of $F$ while meeting none of their interiors.
It is therefore covered by their finitely many boundary segments, which
have empty interior; a finite union of segments cannot contain a nonempty
open set. This is the regular-closedness step used, without being named,
in the paper's Lemma~2 (``covers and hence includes'') and in
Lemma~\ref{lem:finite-neighbourhoods}; both cite this lemma.
\end{proof}

\begin{lemma}[Uniform sector radius and vertex separation in crowns]
\label{lem:sector-radius}
\lean{Tilings.Crowns.uniformSectorRadius}
\uses{def:crown-sector,lem:polygon-local-cone}
There are $r_*>0$ and $\delta_*>0$, depending only on $\mathcal C$, such that
every represented pointed crown $K$ with centre $v$ satisfies
\[
  \bigl(\bigcup K\bigr)\cap B(v,r_*)
       =\Omega(K)\cap B(v,r_*),
\]
and every listed vertex $w\ne v$ of any tile of $K$ satisfies
$\lVert w-v\rVert\ge\delta_*$.  In particular a complete crown covers
$B(v,r_*)$.
\end{lemma}
\begin{proof}
For one crown take the minimum of the finitely many radii in
Lemma~\ref{lem:polygon-local-cone}, one for each incident tile.  This gives
the displayed equality.  Take a positive minimum of these radii over the
finite catalogue, capping them by $1$ if desired.  A largest admissible
radius need not be attained and is not used.

The collection of positive distances from the distinguished centre of a
catalogue crown to a listed vertex of one of its tiles is finite.  Its
minimum is positive; if the collection were empty, use $\delta_*=1$.
Translations preserve both conclusions.  Collinear listed vertices can be
included in this minimum.  A complete crown has cone equal to the whole
plane, so its sector is the ball.
\end{proof}

\begin{definition}[Geometric macro boundary]
\label{def:macro-boundary}
\lean{Tilings.Crowns.geometricMacroBoundary}
\uses{def:polygon,def:protoset,def:substitution}
Let $P$ be the finite prototile family and let $\theta>1$ be the expansion
factor.  A finite patch $M$ considered below has pairwise disjoint tile
interiors and union
\[
 R=u+\theta^j A,\qquad A\in P,\quad j\in\mathbb N.
\]
For geometric arguments merge consecutive collinear edges of the root
polygon $A$.  Their scaled copies are called \emph{geometric macro edges},
and their endpoints \emph{genuine macro corners}.  Every genuine corner
has interior angle different from $\pi$.  Dummy collinear vertices remain
in child-tile and crown data; only the geometric macro boundary is simplified.
\end{definition}

\begin{lemma}[Genuine macro corners are child vertices]
\label{lem:macrocorner}
\lean{Tilings.Crowns.macrocornerIsChildVertex}
\uses{def:tiling,def:macro-boundary,lem:polygon-local-cone}
Every genuine macro corner $q$ is a listed vertex of at least one child
tile of $M$.
\end{lemma}
\begin{proof}
Because $M$ is finite and its carriers are compact, a sufficiently small
ball at $q$ meets only child tiles containing $q$.  If none has a vertex
there, each incident tile has either the whole plane or a halfplane as
its local cone.  The whole-plane case would make $q$ interior to $R$.
Two open halfplanes have nonempty intersection unless they are the two
opposite halfplanes of the same line.  Such intersection would violate
disjoint child interiors.  Thus the local union is one halfplane, or the
whole plane in the opposite-halfplane case.  Neither is the nonstraight
wedge at a genuine macro corner.  This contradiction proves the assertion.

A dummy $180$-degree macro vertex need not be a child vertex.  We never
apply this assertion to such a vertex.
\end{proof}

\begin{lemma}[Corrected distance to a macro edge endpoint]
\label{lem:endpoint-gap}
\lean{Tilings.Crowns.incidentMacroEdgeVertexGap}
\uses{def:crown-sector,def:macro-boundary,lem:sector-radius,lem:macrocorner}
Assume all full crowns of $M$ are represented in $\mathcal C$.  Put
$h=\min(r_*,\delta_*)$.  If $q$ is a genuine macro corner and $v$ is a
child-tile vertex in the relative interior of an incident geometric macro
edge, then
\[
  \lVert v-q\rVert\ge h>0.
\]
\end{lemma}
\begin{proof}
Suppose $\lVert v-q\rVert<h$, and let $B$ be a child tile with vertex $v$.
By Lemma~\ref{lem:macrocorner} the crown $K_M(q)$ is represented in
$\mathcal C$.  Near $v$, both $R$ and the tangent wedge of $R$ at $q$
agree with the inward halfplane of the incident macro edge.  Choose this
neighbourhood inside $B(q,r_*)$ and away from the other ray at $q$.
Lemma~\ref{lem:sector-radius} shows that the inward half-neighbourhood is
covered by the tiles of $K_M(q)$.

Regular closedness of $B$ supplies a nonempty open subset of its interior
arbitrarily close to $v$.  It lies in that inward half-neighbourhood.  If
$B$ were not in $K_M(q)$, this open subset would be covered by the finite
union of its other tile carriers.  The finitely many boundary line
segments cannot cover a nonempty plane open set.  A remaining point would
belong to the interiors of $B$ and of a crown tile, a contradiction.
Hence $B\in K_M(q)$.  Its vertex $v\ne q$ is then among the vertices used
to define $\delta_*$, giving $\lVert v-q\rVert\ge\delta_*$, again a
contradiction.

The restriction to an incident macro edge is essential to this argument.
For a vertex on a nonincident macro edge we use the separate root
separation constant below.  Merely belonging to the closed union of a
crown does not imply that every tile containing the point is in the crown.
\end{proof}

\begin{definition}[Positive root separation constants]
\label{def:crown-separation}
\lean{Tilings.Crowns.rootSeparationConstants}
\uses{def:protoset,def:macro-boundary}
Choose $\nu>0$ bounded above by every distance from a genuine root corner
to a root edge other than its two incident edges.  Choose $\lambda>0$
bounded above by every distance between two nonadjacent geometric root
edges.  These are positive finite minima of distances between disjoint
compact sets.  If a relevant collection is empty, use $1$ for that bound.

For each reflex root corner let $\beta\in(0,\pi)$ be its exterior wedge
angle.  Define
\[
  \mu_* =\min\bigl(\{\pi/2\}\cup
          \{\beta:\text{reflex root corners}\}\bigr).
\]
Thus $0<\mu_*\le\pi/2$.  Convex exterior angles greater than $\pi$ and
dummy straight corners are excluded.  Scaling by $\theta^j\ge1$ does
not decrease the distance bounds $\nu$ and $\lambda$.
\end{definition}

\begin{lemma}[Distance to an adjacent ray]
\label{lem:ray-bound}
\lean{Tilings.Crowns.distanceToReflexBoundaryRay}
\uses{def:crown-separation}
Let two rays from $q$ enclose an angle $\beta\in(0,\pi)$, and let $v$ be
on the first ray at distance $t>0$ from $q$.  Its distance to the second
ray is $t\sin\beta$ when $\beta\le\pi/2$, and $t$ when
$\beta\ge\pi/2$.  For a reflex root angle it is therefore at least
$t\sin\mu_*$.
\end{lemma}
\begin{proof}
For unit directions $a,b$ with $a\mathbin{\cdot}b=\cos\beta$, minimize
\[
  \lVert ta-sb\rVert^2=t^2+s^2-2ts\cos\beta
\]
over $s\ge0$.  The minimum occurs at $s=t\cos\beta$ when this is
nonnegative, and at $s=0$ otherwise.  This gives the two formulas.  On
$(0,\pi/2]$ sine is increasing, and for larger angles the ray-distance
formula gives $t\ge t\sin\mu_*$.  Using the whole supporting line in
the obtuse case would give the wrong estimate.
\end{proof}

\begin{lemma}[Uniform relative crown neighbourhood]
\label{lem:relative-crown-ball}
\lean{Tilings.Crowns.uniformRelativeCrownBall}
\uses{def:crown-sector,def:macro-boundary,def:crown-separation,lem:sector-radius,lem:macrocorner,lem:endpoint-gap,lem:ray-bound}
Put
\[
 d=\tfrac12\min\{r_*,\nu,\lambda,h\sin\mu_*\}>0.
\]
For every child vertex $v$ of $M$,
\[
 R\cap B(v,d)\subseteq\Omega(K_M(v))\cap B(v,d)
                    \subseteq\bigcup K_M(v).
\]
\end{lemma}
\begin{proof}
The local tangent cone of $R$ at $v$ equals the union of the child cones
at $v$: intersect the finite patch with a sufficiently small ball and
use Lemma~\ref{lem:polygon-local-cone}.  If $v$ is interior to $R$, that
cone is the whole plane and its crown is complete.

At a genuine macro corner, a ray outside its tangent wedge is outside
$R$ initially.  A first reentry cannot meet either incident ray again,
so lies on a nonincident macro edge at distance at least
$\theta^j\nu\ge\nu>d$.  No outside-wedge point of $R$ is within $d$.

At an interior point of a geometric macro edge the cone is its inward
halfplane.  A ray into the open outward halfplane cannot reenter through
a nonadjacent macro edge before travelling $\theta^j\lambda\ge\lambda$.
An adjacent edge at a convex endpoint is wholly on the inward side, so
cannot cause reentry.  At a reflex endpoint $q$,
Lemma~\ref{lem:endpoint-gap} gives $\lVert v-q\rVert\ge h$, and
Lemma~\ref{lem:ray-bound} gives distance at least $h\sin\mu_*>d$ to
the adjacent edge.  This excludes reentry in all cases.  Finally
$d\le r_*$ and Lemma~\ref{lem:sector-radius} give the last inclusion.
\end{proof}

\begin{lemma}[A uniform ball across a trimmed tile edge]
\label{lem:trimmed-edge-ball}
\lean{Tilings.Crowns.trimmedEdgeTwoTileBall}
\uses{def:polygon,def:protoset,def:tiling,lem:polygon-local-cone,lem:relative-crown-ball}
For the number $d$ above there is $\eta>0$, depending only on $P$ and
$d$, with the following property.  If translated tiles $A,B$ have
disjoint interiors, their edges overlap near $e$, the tiles lie locally
on opposite sides there, and $e$ has distance at least $d/2$ from every
endpoint of those edges, then $B(e,\eta)\subseteq A\cup B$.
\end{lemma}
\begin{proof}
For each prototile edge take its closed trimmed segment of points at
distance at least $d/2$ from both endpoints.  Whenever nonempty, this
compact segment is disjoint from the compact union of all other boundary
edges of that polygon.  Take a positive lower bound $\eta$ for all these
finitely many distances; use $1$ if there are no nonempty trimmed
segments.  Translations preserve the bound.

Inside $B(e,\eta)$ neither tile has another boundary edge.  Its one edge
separates the ball into two connected half-balls, and the local inward
half-ball remains in its interior because no other boundary can be
crossed.  The two opposite inward half-balls, together with their common
line, cover the ball.
\end{proof}

\begin{lemma}[Exterior excursions across nearby macro edges]
\label{lem:exterior-excursion}
\lean{Tilings.Crowns.noShortExteriorExcursion}
\uses{def:macro-boundary,def:crown-separation,lem:macrocorner,lem:relative-crown-ball}
Let $H$ be convex with diameter less than
$\min\{\lambda,d\sin(\mu_*/2)\}$.  If every point of $H$ has distance
at least $d/2$ from every child vertex, then $H\cap R$ is convex.
\end{lemma}
\begin{proof}
Otherwise a segment joining two points of $H\cap R$ has an open exterior
excursion.  Let $E,F$ be its exit and reentry points on the macro boundary.
The finite polygonal boundary gives such endpoints; a segment contained
in the boundary is not an exterior excursion.  Both endpoints are in $H$.
They cannot lie on the same macro edge, since the segment between two
points of one edge lies in $R$.  Nonadjacent edges give
$\lVert E-F\rVert\ge\lambda$, contradicting the diameter bound.

For adjacent edges let $q$ be their genuine common corner.  It is a child
vertex by Lemma~\ref{lem:macrocorner}, so both distances from $E,F$ to
$q$ are at least $d/2$.  The excursion must cross a reflex exterior
wedge: at a convex corner a chord to the adjacent edge goes into the
inward halfplane at its exit point.  Writing $a=\lVert E-q\rVert$,
$b=\lVert F-q\rVert$ and using the reflex exterior angle $\beta$ gives
\[
 \lVert E-F\rVert^2=(a-b)^2+4ab\sin^2(\beta/2)
       \ge d^2\sin^2(\beta/2)
       \ge d^2\sin^2(\mu_*/2).
\]
Here $\beta/2\in(0,\pi/2)$, where sine is increasing.  Again this
contradicts the diameter bound.  Every segment in $H\cap R$ stays in $R$.
\end{proof}

\begin{lemma}[Convexification preserves the diameter bound]
\label{lem:crown-convexification}
\lean{Tilings.Crowns.closedConvexHullDiameter}
\uses{def:polygon}
If $S$ has finite diameter, its closed convex hull $H$ has diameter no
greater than that of $S$.
\end{lemma}
\begin{proof}
For finite convex combinations $x=\sum_i a_i p_i$ and
$y=\sum_j b_j q_j$ of points of $S$, use
$x-y=\sum_{i,j}a_i b_j(p_i-q_j)$ and the triangle inequality.
The coefficients sum to $1$, giving the desired distance bound.
Passing to limits gives the same bound in the closed convex hull.
\end{proof}

\begin{lemma}[Uniform crown cover for a plane tiling]
\label{lem:crown-cover-plane}
\lean{Tilings.Crowns.uniformCrownCoverPlane}
\uses{def:tiling,def:crown-sector,lem:localfinite,lem:sector-radius,lem:trimmed-edge-ball,lem:crown-convexification}
There is $\varepsilon>0$ such that, in every plane tiling by $P$ whose
full pointed crowns are represented in $\mathcal C$, every set $S$ of
diameter less than $\varepsilon$ is contained in one crown.
\end{lemma}
\begin{proof}
Use $\varepsilon<\min\{\eta,d/2\}$ and replace $S$ by its closed convex
hull $H$.  If $H$ comes within $d/2$ of a tile vertex $v$, all of $H$ is
within distance $d$ of $v$.  The complete crown at $v$ covers this ball.

Otherwise $H$ stays at distance at least $d/2$ from every vertex.  If one
tile contains $H$, take any crown containing that tile.  If not, choose
a tile $A$ meeting $H$, a point of $A\cap H$, and a point of $H\setminus A$.
Their segment exits $A$ at a nonvertex point $e$ of an edge.  Near $e$
incident tile carriers are halfplanes.  Disjoint interiors and plane
coverage force a tile $B$ on the opposite side.  The point $e$ is at
distance at least $d/2$ from both tiles' edge endpoints, so
Lemma~\ref{lem:trimmed-edge-ball} covers $H\subset B(e,\eta)$ by $A\cup B$.
Their shared edge segment has an endpoint which is a vertex of at least
one tile.  The full crown there contains both tiles and hence $H$.
An empty $S$ is covered by any crown.
\end{proof}

\begin{lemma}[Uniform relative crown cover for a supertile]
\label{lem:crown-cover-supertile}
\lean{Tilings.Crowns.uniformCrownCoverSupertile}
\uses{def:substitution,def:crown,def:flc,def:crown-sector,def:macro-boundary,lem:relative-crown-ball,lem:trimmed-edge-ball,lem:exterior-excursion,lem:crown-convexification}
There is $\varepsilon>0$ such that every finite supertile patch $M$ above,
whose full pointed crowns are represented in $\mathcal C$, has a crown
covering $S\cap R$ whenever $\operatorname{diam}(S)<\varepsilon$.
One may take
\[
 \varepsilon=\tfrac12\min
   \{\eta,d/2,\lambda,d\sin(\mu_*/2)\}>0.
\]
\end{lemma}
\begin{proof}
Replace $S$ by its closed convex hull $H$ before dividing into cases.
Lemma~\ref{lem:crown-convexification} preserves the strict diameter
bound.  If $H$ comes within $d/2$ of a child vertex $v$, every point of
$H$ is within $d$ of $v$.  Lemma~\ref{lem:relative-crown-ball} therefore
covers $H\cap R$ by $K_M(v)$.

In the remaining case $H$ stays at distance at least $d/2$ from all child
vertices.  Lemma~\ref{lem:exterior-excursion} shows that $W=H\cap R$ is
convex.  If empty, any crown suffices; if contained in one tile, any crown
containing that tile suffices.  Otherwise choose $A$ meeting $W$ and a
segment in $W$ from $A$ to a point outside $A$.  Let $e$ be an exit point.
It is not a child vertex.  A small neighbourhood meets only incident
tiles, each with a halfplane cone.  A segment leaving $A$ while staying
in $R$ requires an incident tile $B$ on the opposite side of the edge.
Indeed, with only the inward halfplane of $A$ present, such an exit is
locally impossible; a segment tangent to the edge remains in $A$ until
an endpoint, which has been excluded.

All relevant endpoints are at least $d/2$ from $e$.  Hence
Lemma~\ref{lem:trimmed-edge-ball} gives $B(e,\eta)\subset A\cup B$.
Since $e\in W$ and $\operatorname{diam}(W)<\varepsilon<\eta$, this union
covers $W$.  An endpoint of the positive-length overlap of the two edges
is a vertex of at least one of the tiles, and its full crown contains
$A\cup B$.  It therefore covers $W$ and the original $S\cap R$.
The overlap need not be an entire edge of both tiles: the argument does
not assume side-to-side matching.
\end{proof}

\begin{lemma}[Pointing a finite crown catalogue]
\label{lem:pointed-flc}
\lean{Tilings.Crowns.pointedFLCOfFLC}
\uses{def:crown,def:flc,lem:localfinite}
If the allowed full crown patches have only finitely many translation
classes, their pointed versions also have only finitely many translation
classes.
\end{lemma}
\begin{proof}
Choose one finite patch for each unpointed translation class.  A crown
centre is a listed vertex of at least one of its tiles, so each chosen
patch has only finitely many possible centres.  Retain all centre choices
which realize an allowed full crown.  This gives a finite pointed
catalogue, including incomplete crowns whenever the original catalogue
includes boundary crowns of finite supertiles.
\end{proof}

\begin{theorem}[Corrected FLC crown-cover statement]
\label{thm:corrected-crown-cover}
\lean{Tilings.Crowns.flcUniformCrownCover}
\uses{def:substitution,def:flc,lem:pointed-flc,lem:crown-cover-plane,lem:crown-cover-supertile}
For a finite polygonal substitution with the article's FLC condition,
using all allowed full crowns, there is a uniform $\varepsilon>0$ satisfying both
the plane and supertile crown-cover conclusions above.
\end{theorem}
\begin{proof}
Lemma~\ref{lem:pointed-flc} gives representatives for the finite translation
classes of allowed pointed crowns.  Every crown of the tilings in the plane conclusion is
allowed by assumption, and every full crown of a supertile is allowed
by the definition of the substitution catalogue, including those on its
boundary.  Apply the preceding two lemmas, decreasing the two positive
constants to their minimum if needed.
\end{proof}

This proves the geometric statement required by the FLC argument.  It
does not address the independent geometric grouping and routing
obligations for quasi-squares, nor the decoration repair in the
composition construction.

\section{A repaired grouping construction and its geometric obligations}

This chapter replaces the arbitrary tile-to-square assignment in
\S\S3.1--3.4 of the paper.  It applies to simple polygonal disks, including
nonconvex polygons.  It does not assume convex prototiles.  Throughout,
collinear original vertices are retained on the boundaries of grouped tiles.
Consequently a grouped side remains a complete original side and scattering
does not introduce a side-to-side defect.

The construction uses one central group and $4N$ annular groups, instead of
the $4(N+1)$ independent small squares in the original drawing.  The four
corner squares are absorbed into adjacent annular groups.  Every new annular
group meets the centre along one connected arc and supplies one print.
There are still exactly $4N$ marked central sides, as required later.

All Lean names below are planned declaration names.  No declaration in this
chapter is claimed to have been checked by Lean.  The two formerly open
geometric estimates are proved below using polygonal Jordan theory,
piecewise-linear annulus routing, and an explicit compactness argument.
Those classical helper lemmas still require Lean implementations; no
geometric certificate is accepted as an assumption of the final theorem.

\begin{definition}[Mesh constants and the finite cellulation]
\label{def:mesh}
\lean{Tilings.Geometry.MeshConstants}
\uses{def:polygon,def:protoset,def:placed,def:tiling,def:sts}
Let $P$ be a finite family of full-dimensional simple polygonal disks.  Choose
$D>0$, $a>0$, $\gamma>0$, and $K_0\in\mathbb N$ such that every original
tile has diameter at most $D$, area at least $a$, every original side has
length at least $\gamma$, and every original tile has at most $K_0$ sides.

A finite STS patch covering a polygonal disk $\Omega$ is regarded as a finite
planar cellulation: original tiles are its closed faces, original sides are
its edges, and original vertices are its vertices.  A face interior is
connected, and distinct face interiors are disjoint.  We write
$B_\infty(r)=[-r,r]^2$ and use Minkowski sums $Q+B_\infty(r)$ for square
buffers.  A Euclidean diameter bound $D$ also bounds the coordinate diameter.
\end{definition}

The constants exist because $P$ is finite.  Neither FLC nor the crown-cover
lemma is used to obtain them.  In particular the grouping construction must
not depend on the final soficity theorem, on the desired routing certificate,
or on the FLC conclusion derived from that certificate.

\begin{lemma}[Regular grouping with an inside anchor]
\label{lem:regular-grouping}
\lean{Tilings.Geometry.regularGrouping}
\uses{def:mesh,def:tiling}
Let a finite STS patch of mesh $D$ cover $\Omega$, and let $Q$ be a square
with $Q+B_\infty(3D)\subset\operatorname{int}\Omega$.  There is a finite
union $L$ of complete original faces such that
\[
 Q\subset L\subset Q+B_\infty(D),
\]
$L$ is a polygonal disk, and the dual graph of its faces is connected.
Every edge of $\partial L$ adjoins on its inside an original face $F$ such
that $\operatorname{int}F\cap\operatorname{int}Q\ne\varnothing$.
\end{lemma}

\begin{proof}
Select the faces whose interiors meet $\operatorname{int}Q$ and let their
union be $U$.  The open square is contained in $\operatorname{int}U$.
Every selected face interior is connected and meets that square; thus
$\operatorname{int}U$ is connected.  The union is finite and regular closed.
Since every selected face has diameter at most $D$,
$U\subset Q+B_\infty(D)$.

Add every bounded connected component of the \emph{open} complement
$\mathbb R^2\setminus U$, taking its closure together with $U$.  A connected
face interior disjoint from $U$ is contained in one complementary component.
The added region therefore consists of complete original faces; no original
face is cut.  There are finitely many such components, since their boundaries
are contained in the finite original skeleton.

The complement outside the convex square $Q+B_\infty(D)$ is connected to
infinity, so every filled component lies inside that square.  This gives the
support bounds.  Filling a complementary component joins its face interiors
to the existing interior across boundary edges.  Hence the filled interior
is connected, and its open complement is the unbounded complementary
component and is connected.

Here is the classical planar fact needed in the formalization.  A finite
regular closed union of polygonal faces, with connected interior and
connected open complement, has a single simple polygonal boundary.  At a
boundary vertex there cannot be alternating blocks of inside and outside
faces: interior arcs joining the two inside blocks, together with a small
local arc, would separate the two outside blocks by the polygonal Jordan
theorem.  The outside could not be connected.  Therefore the boundary graph
has degree two at every boundary vertex.  It is a union of disjoint cycles;
connected interior and connected exterior exclude more than one cycle.
There are no zero-area appendages, since the region is a union of complete
faces.  Thus the boundary is one simple cycle and $L$ is a polygonal disk.

Filling deletes interior boundary edges and creates no new exterior edge.
The inside face along every remaining boundary edge was selected for $U$;
this proves the anchoring statement.  Finally join two face-interior points
by a generic polygonal path in $\operatorname{int}L$.  Perturb it away from
the finite vertices.  Its consecutive faces share complete sides, giving
a walk in the face-dual graph.
\end{proof}

For Lean, the proof should be divided into a finite-face complement lemma,
a polygonal Jordan-boundary lemma, and a generic-path-to-dual-walk lemma.
The Jordan-boundary lemma is a genuine topological result and must be proved
or obtained from a proved polygonal topology API; the definition of a
``group'' must not simply declare the filled union to be a disk.

\begin{lemma}[The dual outside-anchored cycle]
\label{lem:outside-grouping}
\lean{Tilings.Geometry.outsideAnchoredGrouping}
\uses{def:mesh,lem:regular-grouping}
Let $Q_1+B_\infty(10D)\subset\operatorname{int}\Omega$ and put
$K_1=Q_1+B_\infty(3D)$.  There is a polygonal disk $H$, made of complete
original faces, such that
\[
 Q_1+B_\infty(2D)\subset H\subset Q_1+B_\infty(3D).
\]
Every edge of $\partial H$ adjoins on its outside a face whose interior
meets the open exterior of $K_1$.
\end{lemma}

\begin{proof}
Let $E$ be the closed exterior of $\Omega$, together with every original
face whose interior meets the open exterior of $K_1$.  The open exterior
of $K_1$ is contained in $\operatorname{int}E$, and every selected face
interior meets it.  Thus $\operatorname{int}E$ is connected.

No selected face can enter $\operatorname{int}(Q_1+B_\infty(2D))$:
one of its coordinates would have to change by more than its diameter $D$
between that point and its point outside $K_1$.  Let $H$ be the closure of
the component of $\mathbb R^2\setminus E$ containing $Q_1$.  This gives
the lower inclusion.  The open exterior of $K_1$ lies in $E$, giving the
upper inclusion.  As in the preceding lemma, face interiors belong to
single complementary components, so $H$ is a union of complete faces.

Its interior is connected by its definition.  The open complement
$\mathbb R^2\setminus H$ is \emph{not} simply $\operatorname{int}E$: it
also contains the \emph{pockets}, the other components of
$\mathbb R^2\setminus E$, which consist of unselected faces.  Each pocket
borders selected faces across whole edges, so it is joined to
$\operatorname{int}E$ through those edges, and $\operatorname{int}E$ is
connected and disjoint from $H$.  Hence $\mathbb R^2\setminus H$ is
connected and $H$ has no holes.  A pocket face cannot sit across an edge
of $\partial H$, because the far side of such an edge lies in
$\operatorname{int}E$.  With connected interior and connected open
complement, alternating boundary sectors are excluded by the same
polygonal Jordan argument as before.  Hence $H$ is a polygonal disk.  Each exterior
boundary edge adjoins a face selected for $E$, proving the outside anchor.
The margin $10D$ ensures that this face is an original face of the patch
rather than the artificial unbounded exterior region.
\end{proof}

The distinction between inside and outside anchors matters.  A cycle lying
in a thin band around a square may run backwards along the band for a long
distance.  Its mere Hausdorff proximity to the square does not give a
uniform diameter bound for sectors between selected endpoints.

\begin{lemma}[Disjoint skeletal spokes]
\label{lem:annular-spokes}
\lean{Tilings.Geometry.existsAnnularSpokes}
\uses{def:mesh,lem:regular-grouping,lem:outside-grouping}
Let $N\ge2$ and $R$ be sufficiently large compared with $D$.  Take concentric
squares $Q_0,Q_1$ of side lengths $NR,(N+2)R$ respectively, with the margin
required above.  Construct $L,H$ by the preceding lemmas.  The annulus
$H\setminus\operatorname{int}L$ admits $4N$ mutually disjoint simple arcs
in the original skeleton, each joining $\partial L$ to $\partial H$ and
meeting those cycles only at its endpoints.  Their endpoints are distinct
and appear in the same cyclic order on the two cycles.  (The index $j$ of
the $j$th cut disk is matched with the ideal sector $T_j$ only in the next
lemma, through the walls: each wall crosses $\partial L$ exactly once, and
pigeonhole over the ideal wedges fixes the correspondence.)  The components
obtained by cutting along the arcs are $4N$ polygonal disks consisting of
complete original faces.  Each has a connected face-dual graph and shares
one connected nondegenerate arc with $\partial L$.
\end{lemma}

\begin{proof}
Mark the $4N$ positions subdividing the four sides of $Q_0$ into intervals
of length $R$, including the corners.  Draw the radial segments from the
square centre through these positions to the corresponding positions of
$K_1=Q_1+B_\infty(3D)$.  In the square annulus, successive radial segments
have separation at least a universal positive multiple of $R$.  This
follows by considering a straight side and the two sides adjacent to a
corner; the homothety factor is $1+2/N+6D/(NR)$.  Extending each segment
a bounded number of mesh diameters past the two cycles preserves disjoint
tubes when $R/D$ is sufficiently large.

Follow the finite sequence of original faces crossed by one extended
segment.  Consecutive faces share a side or vertex, and every face boundary
is a connected polygonal cycle.  Their boundary graphs form a connected
finite graph, all of whose edges lie within $D$ of the segment.  It meets
both $\partial L$ and $\partial H$.  Select a graph path, erase loops,
and take the part from its last point on $\partial L$ to its first
subsequent point on $\partial H$.  The part lies in the annulus and has
no other point on either cycle.  Planarity of the skeleton ensures that
an edge cannot cross a cycle without meeting it at a vertex.

The separated tubes give disjoint arcs and distinct endpoints.  Disjoint
arcs joining the two boundary components of an annulus have compatible
cyclic endpoint order.  Cutting a polygonal annulus along such $4N$ arcs
gives $4N$ disks, one between each successive pair; this is a direct
polygonal Jordan argument, or the elementary polygonal annulus theorem.
All cuts use original edges, so the disks are unions of whole original
faces.  Their inner boundary arcs have distinct endpoints and positive
length.  Their face-dual graphs are connected by the generic-path argument
used in Lemma~\ref{lem:regular-grouping}.
\end{proof}

\begin{lemma}[Quantitative sector localization and private neighbours]
\label{lem:quantitative-sector}
\lean{Tilings.Geometry.quantitativeAnnularGrouping}
\uses{def:mesh,lem:regular-grouping,lem:outside-grouping,lem:annular-spokes}
There are finite positive constants $C_{\rm sep},C_{\rm loc},C_{\rm len}$
and $C_{\rm err}$, independent of $N$, the substitution power, and the
particular STS patch, with the following property.  For
$R\ge C_{\rm sep}D$, the cycles and spokes above can be selected so that
each annular group $S_j$ lies in the $C_{\rm loc}D$ neighbourhood of the
corresponding ideal radial trapezoid $T_j$.  Along its outer boundary it
has at least
\[
 b(R,D)=\left\lfloor {R\over C_{\rm len}D}\right\rfloor-C_{\rm err}
\]
distinct fine edges whose outside neighbouring original faces meet no
other annular group $S_t$, $t\ne j$.  In particular, after choosing $R$
so that $b(R,D)\ge K_0+2$, every small group has at least $K_0+2$ ordinary
marked inner sides and at least three private neighbour occurrences
relative to every other small group.

There is consequently a constant $C$ independent of $N$ such that
$\operatorname{diam}S_j\le CR$, and
\[
 Q(P,R)=K_0\left\lceil {4C^2R^2\over a}\right\rceil
\]
bounds the number of retained sides of every small group.
\end{lemma}

\begin{proof}
Translate the common square centre to zero.  Put $A=NR/2$,
$h=A+R+3D$, and $\lambda=h/A$.  Let $v_j$ be the successive length-$R$
positions on $\partial Q_0$; the corresponding outer positions are
$\lambda v_j$ on $\partial K_1$.  Assume $R\ge1000D$.

An inner spoke endpoint $u_j$ satisfies
$A\le\|u_j\|_\infty\le A+D$ and lies within Euclidean distance $D$
of the ray through $v_j$.  If $|u_j-tv_j|\le D$, then
$A-D\le tA\le A+2D$.  Since $|v_j|\le\sqrt2A$,
\[
 |u_j-v_j|\le D+2\sqrt2D<4D.
\]
The outer endpoint $w_j$ satisfies
$h-D\le\|w_j\|_\infty\le h$, giving similarly
$|w_j-\lambda v_j|<4D$.

Use the inside-anchored face at $u_j$ to join it through that face's
interior to $p_j\in\operatorname{int}Q_0$.  The collar lies in the
$5D$ ball about $v_j$.  Join $p_j$ radially to zero inside $Q_0$.
Use the outside-anchored face at $w_j$ to reach
$q_j\notin K_1$, in the $5D$ ball about $\lambda v_j$; continue its
radial ray to infinity outside the convex square $K_1$.

These continuations are disjoint.  On a straight side, adjacent ideal
positions have determinant $AR$.  Perturbations of size at most $5D$
change it by at most $10\sqrt2AD+25D^2$, strictly less than $AR$.
The corner cases are rotations of the same estimate.  Applying it also
throughout the linear perturbation preserves cyclic angular order.
The endpoint boxes are separated by at least $R-10D$; the other radial
segments are separated there by at least $R/4$ before perturbation.
Thus a collar cannot meet another continuation.  Remove any self-loop
within its own endpoint box.  We obtain proper simple walls $W_j$ from
zero to infinity, pairwise disjoint except at zero.  In
$H\setminus\operatorname{int}L$ their portions are precisely the
skeletal spokes and stay within $5D$ of the ideal radial rays.

Let $T_j$ be the ideal trapezoid between successive radial segments and
the square boundaries $Q_0,K_1$.  For a point
$x\in\operatorname{int}K_1\setminus Q_0$ outside the two $5D$ radial
tubes, truncate $W_j,W_{j+1}$ on a sufficiently large circle and close
them along that circle.  The resulting polygonal Jordan loop determines
its wedge.  Deform the core portions inside $Q_0$, straighten the
annular portions inside their convex radial tubes, and deform the
exterior portions outside $K_1$.  The three pieces of each wall are glued
at the two points where the wall crosses $\partial L$ and $\partial H$;
the three homotopies are chosen to agree at these transition points, so
they assemble to one homotopy of the whole loop.  The exterior of $K_1$
is not simply connected, and the exterior deformation is legitimate only
because the exterior portion of a wall is a collar of diameter at most
$5D$ followed by a radial ray: it does not wind around $K_1$, so it is
homotopic, within the exterior of $K_1$ and through curves of increasing
infinity norm, to the ideal radial ray of the same angular order.  None
of these homotopies meets $x$.  Winding number is unchanged, so the actual and ideal wedge
labels agree away from the tubes.  This proves localization there.

If $x$ lies in one of the tubes, clip its nearest radial-line point to
the segment between infinity norms $A$ and $h$.  The clipping moves it
by at most $5\sqrt2D$.  Its distance to the corresponding boundary
segment of $T_j$ is therefore less than
$(1+\sqrt2)5D<13D$.  Taking closures gives the safe bound
\[
 S_j\subset T_j+B(0,20D).
\]
This argument does not assume that the square-band cycles have short or
monotone boundary arcs.

Take the middle half $J$ of the outer edge of $T_j$; its length is
$\lambda R/2\ge R/2$.  Its distance to every other ideal trapezoid is
at least $R/(4\sqrt2)$, by the distance to either bordering radial line.
The inward normal from each $y\in J$ meets $\partial H$ within its
first $D$ of length: $H$ contains the square with halfside $h-D$ and
is contained in the square with halfside $h$.  The middle interval is
at least $R/4$ from the outer corners.  Except at finitely many projected
vertices the intersection lies in an original fine edge $e$.

That intersection has distance at least
$R/(4\sqrt2)-21D>D$ from every other $S_t$.  Its outside original
face, of diameter at most $D$, cannot meet another small group.  Thus
$e$ is an ordinary marked private edge of $S_j$.  The projections of
these private edges cover $J$ apart from finitely many points, and each
has projected length at most $D$.  There are consequently at least
$|J|/D\ge R/(2D)$ distinct private edge occurrences.  Repeated outside
neighbour faces do not affect this occurrence count.

We may take $C_{\rm sep}=1000$, $C_{\rm loc}=20$,
$C_{\rm len}=2$, and $C_{\rm err}=0$.  Since
$\lambda\le2+3D/R$, each ideal trapezoid has span at most $5R$ along
its square side and at most $2R$ across it (the roles of the coordinates
swap for the side sectors), hence diameter less than $6R$.
Localization gives $\operatorname{diam}S_j\le7R$.  A square of side
$14R$ contains $S_j$, so disjoint constituent interiors of area at
least $a$ bound their number by $\lceil196R^2/a\rceil$.
Thus $Q(P,R)=K_0\lceil196R^2/a\rceil$ is admissible.  Retaining all
original vertices proves the stated side bound.
\end{proof}

\begin{lemma}[Inner sector order]
\label{lem:inner-sector-order}
\lean{Tilings.Geometry.quantitativeAnnularGrouping_innerOrder}
\uses{def:mesh,lem:regular-grouping,lem:quantitative-sector}
Let $N\ge2$, $R\ge1000D$, and let $L$ be the inside-anchored group of $Q_0$
with $Q_0\subseteq L\subseteq Q_0+B_\infty(D)$, listed counterclockwise
(winding number $2\pi$ about every interior point).  Let
$y_1,\dots,y_m\in\partial L$ with $y_k$ in the $20D$-neighbourhood of the
ideal sector $T_{j_k}$, where $j_1<\dots<j_m$, and with $y_k$ at distance at
least $50D$ from every other ideal sector $T_t$, $t\ne j_k$.  Then
$y_1,\dots,y_m$ appear in this cyclic order along the counterclockwise
boundary parametrization of $L$.
\end{lemma}

\begin{proof}
This is the inside-anchored counterpart of the outer order statement in the
proof of Lemma~\ref{lem:quantitative-sector} (points of $\partial H$ near
the extended radial segments appear in index order along $\partial H$), and
the remark before Lemma~\ref{lem:annular-spokes} explains why anchoring is
needed: a cycle in the band $Q_0+B_\infty(D)\setminus\operatorname{int}Q_0$
may run backwards along the band for a long distance, so Hausdorff proximity
to $Q_0$ alone does not order its points.  Join every $y_k$ through its
inside-anchoring face (a face of $L$ of diameter at most $D$ whose interior
meets $\operatorname{int}Q_0$) to a point $z_k\in\operatorname{int}Q_0$ with
$|z_k-y_k|\le D$, and then radially to the centre $c$.  These paths lie in
$L$, stay within $D$ of the sectors $T_{j_k}$, which are at least $50D$
apart at the points $y_k$, and meet only at $c$; they cut the disk $L$ into
$m$ sectors whose boundary arcs between consecutive points $y_k,y_{k+1}$
contain no other $y_i$.  A boundary parametrization visiting the points out
of sector order would have to run through the perimeter parameter (the
radial projection to $\partial Q_0$) backwards by at least one sector gap,
which the winding separators of the outer order argument, transplanted to
the inner band with the radial escape paths above in place of the
outside-anchored escape faces, exclude; a parametrization running clockwise
contradicts the winding number $2\pi$ of the counterclockwise listing.  The
constants are those of the Lean statement: $20D$ localization, $50D$
separation, $R\ge1000D$.
\end{proof}

Choose $R$ using that lemma before choosing $N$.  Then choose $N\ge2$ with
$4N>3Q(P,R)$.  Only afterwards choose $k$ so that every $\tau^k$-macrotile
contains the required squares and margin.  Each original polygon has a
chosen interior ball of positive radius; scaling that ball suffices, even
for a nonconvex polygon.  There is no circular dependence of $Q$ on $N$.

Apply the grouping independently to the finitely many fixed patches
$\tau^k\alpha$, $\alpha\in P$.  Their finitely many resulting labelled
groups form $P'$.  This argument needs no claim that arbitrary patches have
only finitely many geometric neighbourhoods.  Such a claim is needed later
for FLC clones and must follow separately from
Lemma~\ref{lem:crown-cover-supertile}, whose proof is independent of routing.

\begin{definition}[Finite geometric macrocertificate]
\label{def:certificate}
\lean{Tilings.Geometry.MacroCertificate}
\uses{def:polygon,def:protoset,def:placed,def:substitution,def:sts,def:mesh}
A macrocertificate consists of finite form and type sets, a finite
side-to-side polygonal template for each form, and the following explicit
finite data and properties.
\begin{enumerate}
\item[(G1)] Every template covers the scaled parent disk, has connected
full inner-side dual graph, and has at least three children.  Every type
has an inner side.  Type tags of different templates are disjoint.
Each internal template edge has its own unique interface identifier,
carried by both incident owners; an inner side is carried by exactly two
children of its template (the side and its partner occurrence).  It records the template, the two type
tags, their incident declared-side slots, and their relative placement;
different interfaces between the same two types remain distinguishable.
In the decorated implementation the identity is
stored in the first component except on a print, where it is stored in
the fifth component.  A special side which is not a print still uses the
first component; specialness alone does not move identity to the fifth.
\item[(G2)] The boundary belt graph is connected.  Its vertices are
\emph{all} types having a vertex on the macroboundary, including types
touching it only at a vertex.  Its edges are inner sides having an endpoint
on the macroboundary.  Each boundary-belt type has such an edge.  The
central child has neither outer nor boundary sides.
\item[(G3)] For each parent declared side, take the ordered list of whole
fine outer sides contained strictly inside its scaled side interval.
At most two fine sides are omitted for crossing its endpoints.  There
are five distinct ports in this list, chosen by a fixed interior-quantile
convention.  Ranks and port numbers are counted in one global positive
direction fixed for each unoriented supporting-line direction, never in
the boundary traversal order of the owner: the quantiles
$\lceil iL/6\rceil$ are not symmetric under reversal (for $L=12$ they are
$2,4,6,8,10$ one way and $11,9,7,5,3$ the other), and the properness
lemma needs port $i$ of a side to face port $i$ of the coincident
opposing side.  Since only translations act, both owners of a coincident
side compute the same direction.  Full matching of a parent declared
interval and its fine tiling then induces identical contained lists and
port numbers.  A parent
dummy endpoint may lie inside a fine child side: no nested subdivision
or hereditary-cut assumption is made.  The finite template also records
the interval intersections with each parent declared side and with the
declared boundary partitions at all required intermediate levels.
Their exact endpoint/frame encoding and matching implications are proved
in the boundary-transport chapter; scalar ordinal ranks alone do not
establish them for arbitrary collinear dummy vertices.
\item[(G4)] Marked parent sides have distinct central prints in cyclic
order.  Cyclic order is the index order of the listed sides, which agrees
with the planar order only when all listed polygons are traversed in the
same rotational sense; the certificate therefore lists \emph{every} form
of its alphabet, the original singleton forms included, counterclockwise
(an original form whose given listing is clockwise is relisted; this
changes side indices but no carrier, side set or scattering).  For every marked side $d$ there are four paths $P_i(d)$ from port
$i$ to its print.  Every side occurrence is used by at most one path,
except the terminal print in either of its two incident owners.  Thus
the central and small owners of that same print are both exempt at the
shared terminal side.  Every path uses marked sides only.  \emph{A path is a sequence of geometric sides, so a genuine path contains both occurrences (both owners) of every inner side it crosses}, the terminal print excepted: the path records the print in its small owner only.  Without this clause the third cell $A_i(d)$ would sit on one occurrence of a crossed side and $0$ on its partner, and the lifting theorem fails for the same reason as below (found by Aristotle, 12 October 2026).  \emph{Every
side of a genuine path other than its initial port is an inner side of the
template}: a path never passes through an outer side of another macroside.
Without this clause the lifting theorem fails, because a path crossing an
outer side $b$ of an adjacent macroside would put a nonzero parent signal
into the third cell of $b$ while the coincident side of the neighbouring
block, lying on no path, carries zero.  The only central side on a genuine
path is its final print.  In its small owner the print is the unique
special side of that owner (the one side of a small type whose cells are
copied from the parent; this is the exception allowed in (G6)).  Its preceding
side is an ordinary side of the small child sharing that print.  The
fifth path and all five paths of unmarked parent sides are singleton
ports, disjoint from genuine paths.  This is a finite routing certificate.
\item[(G5)] There is a unique large central child and all other forms are
original or small.  The central-form map fixes every retained large form.
For a large parent the print map is the identity on its marked sides.
There are $L=4N$ marked central sides, at most $Q$ marked sides on a
noncentral form, and $L>3Q$.  Every original type has a marked inner side.
\item[(G6)] Noncentral fingerprints distinguish types.  For an original
type its fingerprint is the multiset of type components of its ordinary
inner-side identity pairs.  For a small type use its ordinary marked-side
identity pairs and allow at most two further type components from its
unique special side, which is an inner side (the side shared with the
central child).  The possible fingerprint multisets for two
distinct noncentral types are disjoint.  This is a finite separating
property, not an assertion that a composed decorated tile already exists.
\item[(G7)] Marking discipline.  Every outer side, every boundary side
(inner side with an endpoint on the macroboundary), every print, and every
side occurring on a genuine path is marked.  Consequently an unmarked side
is an inner, nonprint, nonboundary side lying on no path; in the
construction these are exactly the sides shared by the central child and
its ring other than the prints.  This condition is what makes the
``local cells'' of Definition~\ref{def:control/typed-decoration} genuinely
type-determined: without it an unmarked boundary side, whose fourth cell
is the parent header, or an unmarked path side, whose third cell is a
parent signal, would carry a parent-dependent value declared local, which
is the error type of the source article.  The paper's own instantiation
satisfies (G7), since it marks every side not belonging to a large
quasi-square.
\end{enumerate}
\end{definition}

The certificate does \emph{not} assume the final soficity statement,
existence of a coloured composition, or a matching-rule system with the
desired language.  These are the conclusions proved from the finite data
in the combinatorial chapter.  Conversely, the final geometric theorem
must construct the certificate from the hypotheses on $\tau$; accepting
the certificate as an unproved input would only prove a conditional core.

\begin{lemma}[A finite sufficient fingerprint criterion]
\label{lem:fingerprint}
\lean{Tilings.Geometry.fingerprintSeparation}
\uses{def:certificate,lem:quantitative-sector}
Suppose (G1) holds, every original form has at most $K_0$ sides, every small
type has at least $K_0+2$ ordinary marked inner sides and one special side,
and, for each pair of distinct small types in one template, the first has
at least three ordinary marked-side occurrences adjacent to original
types which are not ordinary marked-side neighbours of the second.
Then (G6) holds.
\end{lemma}

\begin{proof}
An original fingerprint has at most $2K_0$ components.  A small fingerprint
has at least $2(K_0+2)$ components before its special contribution.  Thus
an original and a small type cannot have equal fingerprints.

Ordinary components of types in different templates lie in disjoint tag
sets.  Every original fingerprint is nonempty.  A small type has more
than two ordinary components, whereas another template can supply at most
two such foreign components through its special side.  This separates
types in different templates.

Consider two original types $u,v$ in one template.  In the fingerprint of
$u$, the multiplicity of $u$ is its inner-side degree, and the multiplicity
of $v$ is the number of its sides adjacent to $v$.  Equal fingerprints
force both degrees to equal the number of $u$--$v$ sides.  All their inner
sides therefore join them to one another.  They form a two-vertex connected
component, contradicting (G1).

For two distinct small types in one template, count the three private
original-neighbour occurrences of the first.  None occurs in the ordinary
fingerprint of the second.  Its special side can add at most two type
components, so it cannot supply all three.  Repeated occurrences of the
same private neighbour are allowed: the argument is about multiset
multiplicity, not distinct neighbour names.  The fingerprints differ.
\end{proof}

\begin{lemma}[Wide routing after a sufficiently large power]
\label{lem:routing}
\lean{Tilings.Geometry.existsRoutingAfterPower}
\uses{def:mesh,def:certificate,lem:quantitative-sector,lem:finite-idempotent}
For the finite grouped family $P'$ constructed above, extend the grouped
substitution $\sigma$ to every group by scattering its original
constituents and applying $\sigma$ to them.  Choose central children as
in the paper and let $c:P'\to P'$ be their form map.  There is an
arbitrarily large power exponent $l$ with $c^{2l}=c^l$ such that, after
retaining the chosen central child and its annular small neighbours and
scattering all other groups in each $\sigma^l$-template, the finite data
required by (G4) exist; the print of every marked parent side is one of the
$4N$ designated marked sides of the central form (the side it shares with one
annular piece), distinct prints on distinct pieces, in cyclic and planar
matching order with the marked parent sides, and the identity on a large
parent.  The setup also records the sector localization of the annular
pieces (Lemma~\ref{lem:quantitative-sector}) and the piece of every print;
the pieces of the prints of the marked sides are pairwise distinct and in
the cyclic order of the marked sides (Lemma~\ref{lem:inner-sector-order},
applied in the construction of Lemma~\ref{lem:certificate-exists}).  The
input is the annular structure of the retained
templates from Lemma~\ref{lem:quantitative-sector}: the central child is the
inside-anchored group $L$, the small neighbours are the $4N$ annular groups
between $L$ and the outside-anchored group $H$ (so every private outer edge
adjoins a face reaching the exterior of $K_1$), $R\ge10000(K_0+3)D$, and
every annular group has four private outer edges in separated windows,
pairwise at distance at least $R/10$; all grouped disks are listed
counterclockwise (as in the construction of
Lemma~\ref{lem:certificate-exists}), so that the polygon loop of a parent
is the positively oriented boundary parametrization used in the
matching-order clauses.  Without the anchoring, the private
faces of a small neighbour could lie in cavities closed off by fewer than
four sides, and the four disjoint paths cannot exist.
\end{lemma}

\begin{proof}
The finite idempotence lemma permits arbitrarily large suitable $l$.
The final chosen centre lies inside the scaled first-generation central
child.  That first-generation child is strictly inside its parent
template.  The finitely many parent templates have a positive minimum
clearance.  After normalizing a $\sigma^l$-template by its expansion
$s=\theta^{kl}$, its central positions consequently lie in a fixed compact
set $K_A\subset\operatorname{int}A$, uniformly over $l$.
For the finite parent catalogue put
\[
 \rho=\min_A\operatorname{dist}(K_A,\partial A)>0.
\]
This handles the moving centre without assuming convergence or convexity.

Outside the retained central annulus every child is an original tile,
so the mesh remains $D$ and every fine outer side has length in
$[\gamma,D]$.  A scaled parent declared side of length $s\ell$ contains
an ordered list of whole fine sides of cardinality
$L\ge s\ell/D-2$: only the at most two sides crossing its endpoints
are omitted.  Choose five ports by the sixth-quantiles
$\lceil iL/6\rceil$, $i=1,\ldots,5$, for sufficiently large $L$.
Their separation from one another and from the ends is bounded below by
\[
 \gamma\bigl(\lfloor L/6\rfloor-1\bigr)-D.
\]
After normalization this gives a positive lower bound depending only on
the finite parent catalogue.  Closed portions of different sides away
from their endpoints also have positive mutual Euclidean separation:
they are disjoint compact subsets of a simple polygonal boundary.
Thus all realized genuine port tuples belong to a compact set of ordered
boundary configurations with uniform separation, including separation
from every other designated or singleton port.  No dummy endpoint is
assumed to be a fine child vertex.

In each selected small group's middle outer interval, choose four
separated windows, for example at relative parameters
$0.30,0.43,0.57,0.70$, and use the projection argument of
Lemma~\ref{lem:quantitative-sector} to choose a private fine edge in
each window.  Let $p_i$ be its midpoint.  Its outside-anchored original
face contains a point $q_i$ outside $K_1$; join $p_i$ to $q_i$ through
that face.  This collar has diameter at most $D$.  The windows are
separated by a fixed multiple of $R$, so the four faces are distinct.
Taking $R\ge10000(K_0+3)D$ makes all these collars and radial rays from
$q_i$ to infinity separated by more than $10D$ outside $K_1$.
The rays remain outside $H$.  These choices are fixed for the finitely
many central annulus catalogues.  Reverse the four entrance labels
where the port numbering is opposite to counterclockwise boundary order;
the paired lists then have the same cyclic order.

For a normalized centre $z\in K_A$, take the intermediate square
$C_z=z+[-\rho/4,\rho/4]^2$.  Its circumradius is less than $\rho/2$,
so it lies strictly in $A$.  Each fixed radial direction meets its
boundary at a point $b_i(z)$.  Consider the polygonal annulus
$A\setminus\operatorname{int}C_z$.  Its ordered outer port points and
inner points $b_i(z)$ admit disjoint simple polygonal arcs.  To prove this
classical helper, cut the annulus along a polygonal crosscut, apply the
piecewise-linear disk theorem, and route ordered lines in the resulting
strip.  Explicitly: for nested simple closed polygons $C\subset\operatorname{int}A$
there is a continuous map $\Psi\colon[0,1]\times\mathbb R\to A\setminus\operatorname{int}C$,
$1$-periodic in the second coordinate and injective on $[0,1]\times[0,1)$,
sending $\{0\}\times\mathbb R$ onto $\partial A$ and $\{1\}\times\mathbb R$ onto
$\partial C$ as loop parametrizations and $(0,1)\times\mathbb R$ onto
$\operatorname{int}A\setminus C$.  It is obtained from the polygonal disk
theorem (ear decomposition of a simple polygon, gluing of disk models along
a diagonal, Alexander-trick boundary reparametrization) applied to the
annulus cut open along one crosscut, by identifying the two copies of the
crosscut; the strip $[0,1]\times[0,1]$ is the cut-annulus model.  The
entrance points and directions below are not a finite catalogue fixed in
advance: the routing statement is uniform over \emph{all} admissible
entrance data, so the compactification argument must include, besides the
normalized centre and the whole port segments, the radial entrance
directions, and the cyclic order of the entrance pieces must be derived
from the planar order of the prints on the contact arcs (the
\texttt{print\_order} datum), not from their indices.  Concretely, the
setup records the piece of every print, and these pieces are in the cyclic
order of the marked sides (Lemma~\ref{lem:inner-sector-order}); every
entrance point lies within $21D$ of the ideal sector of its piece (the
sector localization of Lemma~\ref{lem:quantitative-sector} plus one face of
diameter at most $D$) and outside $K_1$, and entrance points of distinct
marked sides are at least $R/10$ apart, so the radial directions of the
entrances of distinct marked sides are in the cyclic order of the sides
(entrances of adjacent sectors cannot cross their common ray: two points
outside $K_1$ within $21D$ of that ray and of $\partial K_1$ are within
$100D<R/10$ of each other); the four entrances of one marked side are
relabelled so that their radial order matches the port order.  Short disjoint inward-normal outer collars and outward-radial
inner collars can be prescribed first: remove thin attached fingers
around them, route between their tips in the remaining annulus, and
concatenate.  This does not require $A$ to be star shaped.

For one parameter tuple the compact arcs have positive mutual distance
and positive distance from the annulus boundary off these endpoint
collars.  We now prove local stability of these bounds.  A small movement
of a port along its open side is realized in a local rectangle by a PL
slide $(t,n)\mapsto(t+\delta\psi(t,n),n)$, with $\psi=1$ on its
collar and zero at the rectangle boundary.  If
$|\delta|\operatorname{Lip}\psi<1/2$, this is bi-Lipschitz and
preserves the parent disk and the straight normal collar.  The supports
for different ports are disjoint and avoid the central region.

A small change $z\mapsto z'$ is realized by
\[
 H(x)=x+\phi(x)(z'-z),
\]
where a PL bump $\phi$ is one on $B(z,\rho/2)$ and zero outside
$B(z,3\rho/4)$.  The support lies in $\operatorname{int}A$.  Fix the
normalized collar length $\kappa<\rho/4$ once and for all; then the slide
rectangles at the outer ports and the bump support $B(z,3\rho/4)$ are
disjoint from each other and from the outer collars, so the bump fixes
every outer collar and each slide fixes the inner square.
For $|z'-z|\operatorname{Lip}\phi<1/2$, the lower Lipschitz bound
gives injectivity and the contraction equation
$x=y-\phi(x)(z'-z)$ gives surjectivity.  Thus $H$ is a bi-Lipschitz
homeomorphism, translates $C_z$ and its short radial collars exactly,
and fixes $\partial A$.  Compose these maps and shrink the parameter
neighbourhood so the total distortion is at most two.  The old routes
then give routes for every nearby tuple with at least half the old
clearance and the required straight collars.

The centre and port parameter space is compact and there are only
finitely many parent, annulus, and pairing patterns.  Routes are built,
and the local stability argument is applied, for every parameter in the
\emph{closure} of the set of realized tuples, not only for realized ones;
this is legitimate because the separation bounds defining the parameter
space are closed conditions, and it is needed because a finite subcover
is taken of that compact closure.  The subcover therefore gives a uniform
$c_0>0$ for mutual and off-collar boundary clearance in normalized
coordinates.  Choose $l$ larger, still with
$c^{2l}=c^l$, so that $c_0s>100D$ and the physical intermediate square
contains the fixed central annulus.  Join each inner radial ray at its
intermediate-square endpoint to the outer PL route, whose initial
inner collar has the same radial direction.  The two portions meet
only at that endpoint, since one is inside the square and the other
outside it.  Different radial portions are separated by a multiple of
$R$; different outer portions by a multiple of $s$.  Mixed pairs are
also separated: middle outer portions stay $c_0s$ from the whole square,
and distinct endpoint collars are separated from the other rays by the
finite direction separation scaled by $s$.

The first $O(D)$ part of an outer normal collar crosses only faces in
an $O(D)$ ball about its own port, which contains no other port tile.
Afterwards it is more than $D$ from the parent boundary.  The collar may
cross several ordinary original faces; it need not remain in its own
port face.  It avoids every other designated and singleton port tile,
which is the condition actually required.

Finally perturb the routes away from the finite original vertices and
list their crossed faces.  Consecutive faces share whole sides.  Erase
repeated-face loops while preserving the first port face and final
outside-anchored face.  Begin the resulting side sequence with the
port, end it with the private entry edge, and append the print.  The
entry and print belong to the same small group, hence are consecutive
sides of a path in the paper's exact sense.  Different routes cannot
cross the same original face: their $D$ tubes are disjoint away from
the collars, their inner collar faces are distinct, and their outer
collar balls are separated.  Thus every nonterminal side occurrence
is used at most once.  The shared terminal print is exempt in both
its central and small owners.  No central side is met earlier.
All traversed original sides, small entry sides, and prints are marked.
Every singleton port path remains disjoint from the genuine routes.
\end{proof}

\begin{lemma}[Construction of the geometric macrocertificate]
\label{lem:certificate-exists}
\lean{Tilings.Geometry.existsMacroCertificate}
\uses{def:substitution,def:sts,def:certificate,lem:regular-grouping,lem:outside-grouping,lem:annular-spokes,lem:quantitative-sector,lem:fingerprint,lem:routing,lem:finite-idempotent,lem:inner-sector-order}
For every STS polygonal substitution $\tau$, there are a finite grouped
family $P'$, a scattering map $\Delta$, an integer $m>0$, and a grouped
substitution $\sigma'$ satisfying
\[
 \Delta\sigma'\alpha=\tau^m\Delta\alpha
 \quad\hbox{for all }\alpha\in P',
\]
together with a macrocertificate (G1)--(G7).  The grouped family contains
every original form as a singleton group and every group is a group of tiles
of a supertile (the hypotheses of
Lemma~\ref{lem:grouping-lift}), and when $\tau$ has hereditary listed
sides so has $\sigma'$.
\end{lemma}

\begin{proof}
Choose $R$, then $Q$, then $N$, then $k$ in the noncircular order described
above.  Group each fixed $\tau^k\alpha$ by the inside/outside cycles and
the spokes.  Keep the finite labelled groups in $P'$.  Their supports are
disks and their constituent dual graphs are connected.  Scattering a group
and applying the grouped substitution to each constituent defines $\sigma$
and gives $\Delta\sigma=\tau^k\Delta$.

\emph{Step 3 (the routing setup).}  For every grouped form $\alpha$ the
template $\sigma\alpha$ is the regrouping of $\tau^k\Delta\alpha$: each
constituent $q\in\Delta\alpha$ contributes the translated groups of its
supertile $\tau^kq$ (its $L_q$, its $4N$ annular pieces, singletons).  The
\emph{central child} $c(\alpha)$ of $\sigma\alpha$ is the group $L_q$ of the
constituent $q$ whose centre $z_q$ is nearest to the centre of $\alpha$; thus
$c$ maps grouped forms to large catalogue forms $L_p$.  For a power $l$ the
retained template of $\alpha$ keeps, inside $\sigma^l\alpha$, only the
central chain $c^l(\alpha)$ (the central child of the central child, $l$
times) and its $4N$ annular neighbours, and scatters every other group into
original tiles.  Its annulus is a translate of the catalogue annulus of one
original supertile, so the quantitative data ($N$, $R$, the four separated
windows of every piece, the private-edge count) are those of the grouping
step.  The clearance of the central child grows geometrically along the
central chain: at each step the chosen centre lies in the square of
Lemma~\ref{lem:quantitative-sector} with its margin, strictly inside the
scaled parent, so after normalization by $\theta^{kl}$ the central child
stays at a positive distance $\rho$ from the boundary of the retained region,
uniformly in $l$ and $\alpha$ (a geometric series with ratio $\theta^{-k}$).
The marked sides are fixed by the marking rule: every side of a form that is
not a central image is marked; a large form $L_p$ has exactly $4N$ marked
sides, the side it shares with each of the $4N$ annular pieces (one per
piece; these are whole sides of $L_p$ by the whole-side property of groups).
The print slots of $\alpha$ send its marked sides injectively to contact
sides of the central child, preserving cyclic order and in planar matching
order with the marked sides (both disks are listed counterclockwise, and an
order-preserving injection between two cyclically ordered boundary families
is in matching cyclic order in the annulus between them); this is possible
because a non-large form has at most $Q<4N/3$ sides while the centre has $4N$
contact sides, and for a large parent $\alpha=c^l\beta$ the print map is the
identity (its $4N$ marked sides are its own contact sides).  The prints
are chosen sector by sector: for the marked sides $a_1<\dots<a_m$ in index
order take distinct pieces $j_1<\dots<j_m$ (all $4N$ pieces for a large
parent) and for each $j_k$ the contact side of the central child met by the
ray from the centre through the middle of the sector $T_{j_k}$; by the
sector localization the face across that side lies in the piece $j_k$, so
the piece map preserves cyclic order by construction, and by
Lemma~\ref{lem:inner-sector-order} these sides occur counterclockwise in
the order of the sectors, so the print map preserves cyclic order as well;
the sector localization itself is retained from
Lemma~\ref{lem:quantitative-sector}.

\emph{Step 3$'$ (the routed power).}  Lemma~\ref{lem:routing} applied to this
setup gives $l>0$ with $c^{2l}=c^l$ (Lemma~\ref{lem:finite-idempotent}) and
routing data for every retained template at the power $l$.  Keep only central
large forms in $\operatorname{im}c^l$.  For a retained large form
$\alpha=c^l\beta$, $c^l\alpha=c^{2l}\beta=c^l\beta=\alpha$.  Thus the central
form is fixed.  Put $m=kl$ and call the resulting substitution $\sigma'$.
Scattering commutes with its decomposition, proving the displayed
compatibility.

\emph{Step 4 (the finite certificate data).}  The forms of the certificate
are the retained catalogue tiles (the central chain forms, the annular
pieces, and the original singletons), each with its disk relisted
counterclockwise and translated to the origin.  The types are the positioned
children of the retained templates; the owner of a type is its template and
its anchor is its position.  A type is large, small or original according as
its form is a central image, an annular piece or an original singleton.  The
partner of an inner side occurrence is the other occurrence of the same
geometric side in the same template (templates are side-to-side, so an inner
side has exactly two owners, Step 5); the interface identifier of an inner
side is its unordered pair of occurrences, with the two type tags recorded in
a fixed order.  Marking follows the marking rule of Step 3 through the parent
sides.  The special side of a small type is its designated contact side with
the central child: the print when the type receives one (prints land on
distinct pieces), an arbitrary fixed contact side otherwise.  The print of a
marked parent side is the routing's print slot.  The outer list of a parent
declared side is the list of whole outer fine sides of the template contained
in the scaled side, without repetition, in the lexicographically positive
direction of the supporting line.  The genuine paths are the occurrence
chains of the routed side lists: for a routed side list $e_0$ (the port),
$e_1,\dots,e_r$ (the print) with crossed children $X_0,\dots,X_{r-1}$ the
chain is $(X_0,e_0),(X_0,e_1),(X_1,e_1),(X_1,e_2),\dots,(X_{r-1},e_r)$: both
occurrences of every crossed inner side, the print in its small owner only;
its side projection with consecutive repetitions removed is the routed list.
The fifth path and all paths of unmarked sides are singleton ports.  The
chains inherit from the routing data that every side after the port is
inner, that every side is marked, that no central side occurs before the
print, that the side before the print is an ordinary side of the small
owner, that no occurrence repeats, and that distinct paths share only the
terminal print in either owner.  The central child of a template is owned
by that template; the mixed form of a type's child is the catalogue form of
its own form; each mixed form has one catalogue entry; and the relabelling
between the counterclockwise listing of a certificate disk and the listing
of its grouped disk is a bijection of side indices.  The
constants are $N$ and $Q$ as chosen, and the rank bound is the largest outer
list length.  This data represents $\tau$ at the power $m$: $\Delta\sigma'=
\tau^{kl}\Delta$ because scattering commutes with the retained/scattered rule;
the dissections are side-to-side with connected duals, whole sides and
retained vertices because they are catalogue groups; every original form is a
singleton; every group lies in a supertile; heredity transfers.  The retained
templates are side-to-side because scattering a group of a side-to-side patch
along whole sides preserves the side-to-side property.

\emph{Step 5 (G1).}  Disk cellulations with at least the centre and its ring
have connected inner-side duals and at least three children.  Every child has
an inner side: a child whose sides are all outer is the whole template.  An
inner side has exactly two owners, because the template is side-to-side and
three polygons with disjoint interiors cannot share a whole segment.  The
partner and identifier clauses are the definitions of Step 4.

\emph{Step 6 (G2).}  Around the macroboundary, consecutive boundary faces are
connected through the finite fan at each boundary vertex.  Its connecting
edges are precisely inner sides with an endpoint on that boundary.  Faces
touching the boundary only at a vertex must be included in this fan.  Every
boundary type has such an inner edge because it is not the whole macrotile.
The retained centre is strictly interior, by the clearance of Step 3.

\emph{Step 7 (G3).}  The lexicographic order is a global positive
orientation.  The finite boundary interval intersections give the outer
lists; at most two fine sides are omitted for crossing the endpoints of a
scaled side (Lemma~\ref{lem:whole-ports}); each list has at least five
entries because the five ports of the routing are distinct entries of it
(they are the quantiles of this very list); the lists are bounded by the rank
bound.  Ranks and ports are numbered in the global supporting-line direction,
as (G3) requires.  No nested-cut assumption is used.  The exact matching
implications of the endpoint/frame schedules are proved separately in the
boundary-transport chapter.

\emph{Step 8 (G4).}  The print map is the routing's print slot map: injective
and cyclic-order preserving on the marked sides, with inner prints whose other
owner is a small child, where the print is the special side.  The occurrence
chains of Step 4 inherit from the routing that every side after the port is
inner, that every side is marked (sides of non-large children and prints),
that no central side occurs before the print, that the side before the print
is an ordinary side of the small owner (the entry edge and the print belong
to the same small group), that no occurrence repeats, and that distinct paths
share only the terminal print in either owner (their crossed faces are
distinct).

\emph{Step 9 (G5).}  The central child belongs to its template and is its
unique large child.  For a retained large form $\alpha=c^l\beta$,
$c^l\alpha=\alpha$, so the central form is fixed and the print map is the
identity on its marked sides (the central child has the parent's own disk).
A large form has exactly $4N$ marked sides (Step 3); a non-large form has at
most $Q$ marked sides, because all its sides are marked and an annular piece
has at most $Q$ sides (Lemma~\ref{lem:quantitative-sector}) while an original
tile has at most $K_0\le Q$ sides; $4N>3Q$ by the choice of $N$.  Every
original type has a marked inner side: all its sides are marked and it has an
inner side (Step 5).

\emph{Step 10 (G6).}  The special sides are those of Step 4: inner, shared
with the centre, unique per small type.  Every small type has at least
$K_0+2$ private outer edges (Lemma~\ref{lem:quantitative-sector}), so the
private-neighbour estimate and Lemma~\ref{lem:fingerprint} give the
disjointness of the possible fingerprint multisets of distinct noncentral
types.

\emph{Step 11 (G7).}  The marking rule ``every side not belonging to the large
central group, plus the prints'' gives (G7): outer and boundary sides belong
to non-large children (the centre is strictly interior), prints are marked,
and every occurrence on a genuine path lies in a non-large child, since no
central side is met before the print.

This constructs the finite geometric certificate.  Neither a coloured
composition theorem nor the final soficity conclusion has been assumed.
All Lean implementations, including the classical polygonal topology
helpers used in this chapter, remain unverified.
\end{proof}

\paragraph{Why a citation to an earlier matching-rule theorem is insufficient.}
Fernique--Ollinger's theorem assumes a \emph{good} combinatorial
substitution; its connecting condition is the very information-routing
condition at issue here.  The implication from arbitrary planar STS data
to that condition must still be proved.  Goodman--Strauss assumes
hereditary sides and sibling-side-to-side conditions, and directly treats
substitution tilings.  Using it as an alternate route requires verification
of those exact hypotheses and a proved theorem supporting the hierarchical
variant as well.  Neither a bibliographic citation nor an unproved Lean
axiom supplies these constructions or their Lean implementations.

\paragraph{Triangulation is not an unproved refinement reduction.}
The suggested polygonal triangulation is only an auxiliary tool in the
continuous wide-routing proof.  Replacing prototiles by triangles would
require a separate finite substitution-compatible refinement and a proof
of projection equivalence for both substitution and hierarchical tilings.
No such reduction is assumed here.

% A chapter fragment. The main document supplies theorem environments and
% the leanblueprint macros \lean, \leanok, and \uses.
% Only the two HeaderKernel nodes below are marked leanok.

\section{A boundary-header repair of coloured composition}

This chapter proves the hereditary-side core composition statement. The
flexible zero/rank definition requires hereditary listed sides, as defined
in Definition~\ref{def:hereditary}: (G3) allows a fine outer side to cross
a dummy parent endpoint, and for such a side the clause ``$b\subset\sigma'd$''
of Definition~\ref{def:control/substitution} is undefined. Accordingly
\emph{every theorem and lemma of this chapter carries the standing
hypothesis that the macrotemplate system has hereditary listed sides}, so
that each fine outer side lies in a unique scaled parent listed side; the
hypothesis is repeated in the statements of Theorems~\ref{thm:compose},
\ref{lem:lift-hierarchy} and~\ref{thm:comp/soundness} and in their Lean
targets. Original nonhereditary dummy cuts are
handled separately by the strict schedule proof in the final chapter.
This chapter proves a corrected composition statement for an already
constructed macrotemplate system. It is conditional on the finite geometric,
routing, and separation certificate of Definition~\ref{def:certificate}.
It does not establish that certificate for every STS substitution. In
particular, the arbitrary assignment of tiles to grid squares in the source
article does not establish connectedness or the required routing data.
Every internal interface has a distinct finite ID naming its exact two
owner-side slots and relative displacement. Identity data are triples
$(s,r,e)$; fingerprints project them to the ordered type pair $(s,r)$,
so all occurrence counts are unchanged.
All declarations in the namespace \texttt{Tilings.Planned} below are proposed
formalization targets, with mathematical proofs given here; they have not
been compiled. Only the two explicitly marked nodes in
\texttt{Tilings.HeaderKernel} refer to the compiled Lean kernel.

The certificate interfaces used here must include the following conventions.
Outer sides are marked. A child's type determines its owner template and
therefore the form of its parent. Print maps depend only on that form, and
copy-source maps depend only on the child position and the parent form.
The boundary belt is connected, the central tile has no outer or boundary
side, and identity data determine rigid relative placements. All polygon
patches have disjoint interiors and cover their parent polygon. These
conventions are essential uses of the certificate, rather than conclusions
derived from normality alone.

\begin{definition}[Typed decorations and normality]
\label{def:control/typed-decoration}
\lean{Tilings.Planned.typedDecoration}
\uses{def:certificate,def:substitution}
Let $F$ be the finite set of forms and $T$ the finite set of positioned child
types. Write $\operatorname{form}:T\to F$, and let $S(t)$ be the finite
indexed side set of $\operatorname{form}(t)$. The owner of a child type is
its parent form, not necessarily its own form. Side identifiers are retained
even when consecutive polygonal sides are collinear.

An index belongs to a finite tagged alphabet $U$. Its tags distinguish zero,
positive alignment ranks, type labels, exact interface identities, the symbol
$\bot$, and the headers defined below. A decoration of type $t$ consists
of a five-tuple $c_A(d)\in U^5$ on each $d\in S(t)$; its centre stores $t$.
In the hierarchical variant considered here, position $2$ is zero. The
geometric polygon has form $\operatorname{form}(t)$.

The local cells are all cells of an unmarked side, the fifth cell of a special
side, and the first and fifth cells of a marked ordinary side. A decoration
is \emph{normal} when these cells have the values prescribed by its type.
Thus ordinary inner first cells contain the relevant ordered interface identity.
On a special side, the local fifth value is the interface identity if that side
is a print in its owner template, and zero if it is not a print. In particular,
special does not mean print. Call a decoration \emph{identity-correct} when
every inner nonprint side carries its prescribed identity in position $1$
and every print carries it in position $5$. Normality alone does not impose
the first value on a special nonprint side. In the flexible variant of
this chapter an outer fourth cell is not local: an outer side is marked,
and its zero-versus-rank value depends on the parent's mode. In the strict
variant of the final chapter the outer fourth cell is the generated
schedule $W$, which is determined by the type alone and is therefore local
there. The list of local cells relies on the marking discipline (G7): an
unmarked side is inner, nonprint, nonboundary and on no path, so all its
cells are fixed by the type.
\end{definition}

\begin{lemma}[A tile with an outer side has a boundary side]
\label{lem:comp/outer-boundary}
\lean{Tilings.Planned.boundary_side_of_outer_side,Tilings.Planned.Strict.exists_boundary_side_of_outer}
\uses{def:certificate}
In a macrotemplate every child type with an outer side also has a boundary
side, that is, an inner side with an endpoint on the macroboundary.
\end{lemma}
\begin{proof}
By (G1) the type has an inner side, so not all of its sides are outer.
Walk around its boundary polygon from an outer side; the first side that
is not outer starts at the endpoint of the last outer side, which lies on
the macroboundary. That side is inner (the template covers the parent
polygon, so a side not on the macroboundary is shared) and has an endpoint
on the macroboundary. Collinear listed vertices do not affect the walk.
\end{proof}

\begin{definition}[The finite parent header]
\label{def:control/header}
\lean{Tilings.HeaderKernel.makeHeader}
\leanok
\uses{def:control/typed-decoration}
For a type $t$, a Boolean outer-side predicate $o_t:S(t)\to\{0,1\}$,
and a fourth-index vector $a:S(t)\to U$, set
\[
 z_{t,a}(d)=o_t(d)\mathbin{\wedge}[a(d)=0],\qquad
 H(t,a)=\bigl\langle t,z_{t,a}\bigr\rangle
 \in\sum_{t\in T}\bigl(S(t)\to\{0,1\}\bigr).
\]
The Lean definitions are \texttt{Header}, \texttt{zeroMode}, and
\texttt{makeHeader}, with precisely these dependent types and formulas.
The sum is finite when $T$ and each $S(t)$ are finite; the compiled instance
\texttt{headerFintype} uses classical decidable equality and is
\texttt{noncomputable}. For a decorated tile $A$ write $H(A)$ for the header
computed from its type and its fourth vector. This is computed data, not a
new local field fixed by the type.
\end{definition}

\begin{lemma}[Recovery of a header from fourth-port data]
\label{lem:comp/port-recovery}
\lean{Tilings.HeaderKernel.makeHeader_eq_of_port_recovery}
\leanok
\uses{def:control/header}
Fix a type $t$ and fourth vectors $a,b:S(t)\to U$. If a vector
$p:S(t)\to U$ satisfies $p(d)=a(d)$ and $p(d)=b(d)$ for every $d$, then
$H(t,a)=H(t,b)$.
\end{lemma}
\begin{proof}
 \leanok
For each $d$, transitivity gives $a(d)=p(d)=b(d)$. Consequently
$o_t(d)\wedge[a(d)=0]=o_t(d)\wedge[b(d)=0]$. Function extensionality
identifies the two Boolean vectors, and their type components are already
the same. This is the proof implemented in the kernel. If types are initially
identified by an equality $t=s$, their dependent side sets must first be
transported along that equality; the compiled auxiliary lemma
\texttt{makeHeader\_eq\_of\_aligned\_fourth} performs this transport.
\end{proof}

\begin{definition}[The repaired decorated substitution]
\label{def:control/substitution}
\lean{Tilings.Planned.headerSubstitution}
\uses{def:certificate,def:substitution,def:control/header}
The uncoloured substitution $\sigma'$ and its finite templates are unchanged.
For a decorated parent $A$, its repaired decomposition $\widehat\sigma A$
has the prescribed child types. Prints carry the first four cells of the
corresponding parent side, and their fifth cell is the interface identity of the
shared interior side. This prescription applies to \emph{both} owners of a
print. A side which is not a print, including a special nonprint side,
has its prescribed first identity cell, zero second
cell, zero fifth cell, and third cell equal to $A_i(d)$ when it belongs to
the path $P_i(d)$, and zero otherwise. Its fourth cell is
\[
 B_4(b)=
 \begin{cases}
  0,&b\subset\sigma'd\text{ is outer and }z_A(d)=1,\\
  \operatorname{rank}(b),&b\text{ is outer and }z_A(d)=0,\\
  \operatorname{header}(H(A)),&b\text{ is boundary},\\
  \bot,&b\text{ is any remaining inner nonprint side}.
 \end{cases}
\]
The print prescription has precedence. Prints lie strictly inside the
macrotemplate and hence do not conflict with the boundary clause. Each
header tag is distinct from zero. The change from the article is the
boundary clause: it carries the entire parent header, rather than only the
parent type. This introduces finite control data and no new paths.
\end{definition}

\begin{definition}[Copy sources and origins]
\label{def:control/origin}
\lean{Tilings.Planned.copySource}
\uses{def:control/substitution,def:certificate}
Copied cells are positions $1$--$4$ on prints and position $3$ on ordinary
sides assigned to paths. Their source is the corresponding parent cell.
The source is determined by the child position and parent form. All other
cells, including the centre type and the synthesized boundary header, are
origins. Thus a boundary header is \emph{not} falsely treated as a literal
copy of the sixth cell.

Along an ancestor chain, repeatedly follow the source of a copied cell.
If a noncopied cell is eventually reached, it is the origin of the chain.
Otherwise the copied cell has no origin. A print or path may copy a
header-valued fourth cell; its source chain then stops at the synthesized
boundary cell. Values are unchanged along actual copy edges.
\end{definition}

\begin{definition}[Legality and properness]
\label{def:control/legal}
\lean{Tilings.Planned.headerLegal}
\uses{def:control/origin,def:control/substitution,def:tiling}
A decoration $A$ is legal when it has an infinite ancestor chain
$A=A_0,A_1,A_2,\ldots$, with $A_i$ a child of $A_{i+1}$, such that every
copied cell at every $A_i$ with no origin has value zero. All types and
colours lie in fixed finite sets; legality selects a subset of that finite
decorated tile set.

A side is nontrivial if its fourth and fifth cells are not both zero.
A plane tiling is proper if every nontrivial side has a complete, coincident
opposing side with the same full five-tuple. A trivial outer side may have
shifted contacts. Equality of colours is required only for contacts forced
by this rule.
\end{definition}

\begin{lemma}[Elementary legality properties]
\label{lem:comp/legal-properties}
\lean{Tilings.Planned.headerLegal_properties,Tilings.Planned.Strict.strictLegal_properties}
\uses{def:control/legal,def:control/typed-decoration,def:control/substitution}
A legal tile has a legal parent and is normal and identity-correct;
every child of a legal parent
is legal. All boundary fourth cells of a legal tile equal one parent header.
Third signals belonging to the same path are equal through a legal tile.
Every inner side of a legal tile is nontrivial. On an outer side, triviality
is equivalent to a zero fourth cell.
\end{lemma}
\begin{proof}
Delete the first term of a witnessing ancestor chain to obtain legality of
its parent. The child clauses fix every local cell, giving normality, and
supply first identities on all inner nonprint sides and fifth identities
on all prints, giving identity-correctness.
To prove legality of a child, prepend it to a chain for the parent. A copied
cell in the new first tile either reaches an origin there, or reaches the
corresponding parent cell. If the source chain is infinite, the parent's
legality forces its value, and hence the child's value, to be zero. All
requirements farther up the chain are unchanged.

The boundary and path equalities follow by choosing the first legal parent:
its single header supplies every boundary cell, and its single cell $A_i(d)$
supplies all appearances of a path signal in that child. At a print, the
fifth cell is a nonzero interface identity. At a boundary side, the fourth cell
is a nonzero header tag. At any other inner side, the fourth cell is $\bot$.
These cases exhaust inner sides. On outer sides the fifth cell is zero,
so the last assertion follows directly.
\end{proof}

\begin{lemma}[Assembly from active identity sides]
\label{lem:comp/assembly}
\lean{Tilings.Planned.assembleMacrotemplates,Tilings.Planned.Strict.strict_assembleMacrotemplates}
\uses{def:certificate,def:control/typed-decoration,def:control/legal,def:tiling}
A proper plane tiling of normal identity-correct typed tiles in which every
inner side is nontrivial admits a partition into positioned $\sigma'$ macrotemplates.
In particular, every proper legal tiling has such a partition.
\end{lemma}
\begin{proof}
On an inner nonprint side, identity-correctness gives its ordered identity
pair in the first cell. On a print, the fifth cell gives that pair.
Nontriviality forces a complete matching side. The pair identifies the
neighbouring type and the relative placement specified by the certificate.
Starting with one tile, traverse the connected finite inner-side graph of
its owner template. Each graph edge forces its required neighbour. Relative
placements agree because they are the placements of that fixed template;
an existing tile cannot occupy a different positive-area position while
matching the same rigid side placement.

Traversal therefore recovers the full template. Repeated traversal cannot
produce an additional occurrence at the same position, since tile interiors
are disjoint. Two recovered blocks sharing a tile have the same owner type
set and the same rigid placement, and hence are the same block. Every tile
belongs to a recovered block, so the blocks partition the plane. Trivial
outer contacts are not used in this argument. The final assertion follows
from Lemma~\ref{lem:comp/legal-properties}. Notice the hypothesis that inner
sides are active and identities are correct: normality alone would neither
force adjacency on an ordinary marked side whose fourth and fifth cells
were both zero, nor prescribe the first identity on a special nonprint side.
\end{proof}

\begin{lemma}[A common boundary header]
\label{lem:comp/boundary-header}
\lean{Tilings.Planned.boundary_header_constant,Tilings.Planned.Strict.strict_boundary_header_constant}
\uses{def:certificate,lem:comp/legal-properties,def:control/header}
Let $M$ be a proper finite coloured realization of the template $\sigma'\alpha$,
with the prescribed child types. Assume every tile except possibly the
centre is legal. Every boundary fourth cell then carries the same header
$h=(t,v)$, and $\operatorname{form}(t)=\alpha$.
\end{lemma}
\begin{proof}
Every boundary tile is noncentral and legal. Its boundary cells all equal
the header of any chosen legal parent. If two boundary tiles share a boundary
side, that side is nontrivial; properness identifies the two header tags.
Connectedness of the boundary belt propagates equality to every boundary
cell of $M$. Choose one boundary tile $B$ and a legal parent $C$ of it.
The type of $B$ determines the owner form $\alpha$, so $C$ has form $\alpha$.
Its header equals $h$; equality of type components gives $\operatorname{type}(C)=t$.
Thus $t$ has form $\alpha$, as required. Connectedness of this belt is an
explicit certificate obligation, not an inference from a sufficiently
large substitution scale.
\end{proof}

\begin{lemma}[Parent reconstruction and its exact header]
\label{lem:comp/reconstruction}
\lean{Tilings.Planned.reconstructed_header_eq,Tilings.Planned.Strict.strict_reconstructed_header_eq}
\uses{lem:comp/boundary-header,lem:comp/port-recovery,def:control/typed-decoration}
In the situation of Lemma~\ref{lem:comp/boundary-header}, define $D$ to have
type $t$ and, for every parent side $d$ and $1\le i\le5$, set
\[
 D_i(d)=\text{the third cell of port }i\text{ on }\sigma'd.
\]
Then $D$ is normal and identity-correct, $H(D)=h$, and every inner side of
$D$ is nontrivial.
\end{lemma}
\begin{proof}
Fix a local cell $(i,d)$ of $D$. Its port lies in a legal boundary tile $B$,
which has a boundary side by Lemma~\ref{lem:comp/outer-boundary} and hence
a legal parent $C$ with $H(C)=h$.
In particular $C$ has type $t$. The port clause says $D_i(d)=B_3(p_i)=C_i(d)$.
Since $C$ is normal of type $t$, this is the required local value of $D$.
Applying this separately to every local cell proves normality; no equality
of the whole parents $C$ chosen for different ports is needed.

For each $d$, choose its fourth port and a legal parent $C_d$ of that port
tile. We have $H(C_d)=h$ and $D_4(d)=(C_d)_4(d)$. Thus
\[
 v(d)=[d\text{ outer in }t]\wedge[(C_d)_4(d)=0]
     =[d\text{ outer in }t]\wedge[D_4(d)=0]=z_D(d).
\]
The dependent side sets have first been aligned using the established
equality of parent types. Extensionality gives $H(D)=h$. The kernel recovery
lemma is the algebraic component of this argument; obtaining the legal
parents and their common type remains part of this mathematical lemma.

For an inner nonprint side $d$, use its first port and a legal parent of
that port tile. Its type is $t$, so the child clauses force its first cell
to be the prescribed identity of $d$. This port recovers exactly $D_1(d)$.
For a print, the fifth cell of $D$ is its nonzero prescribed identity by
normality. Thus $D$ is identity-correct. For an inner nonprint side, the
legal parent $C_d$ chosen at its fourth port has a nonzero fourth cell: it
is a header or $\bot$. Hence $D_4(d)\ne0$. For a print, the normal fifth
identity cell is nonzero. Therefore every inner side of $D$ is nontrivial,
even before legality of $D$ has been established. This argument explicitly
covers special nonprints, whose fifth cell is zero.
\end{proof}

\begin{lemma}[Signals reach the prints]
\label{lem:comp/signals}
\lean{Tilings.Planned.path_and_print_recovery,Tilings.Planned.Strict.strict_path_and_print_recovery}
\uses{def:certificate,lem:comp/legal-properties,lem:comp/reconstruction}
Every nonprint third signal on $P_i(d)$ in $M$ equals $D_i(d)$.
The noncentral owner of the print of $d$ has first four cells
$(D_1(d),\ldots,D_4(d))$ on that print.
\end{lemma}
\begin{proof}
Begin at port $i$, whose third cell is $D_i(d)$ by construction.
Consecutive nonprint sides of the path belong to the same noncentral legal
tile, or are common sides of two such tiles. In the first case, their
signals agree because both copy the same parent cell. In the second case,
the common side is inner and nontrivial, so properness gives equality of
its signals. Induction along the finite path proves the first assertion.

Let $B$ be the small noncentral tile sharing the final print, and choose a
legal parent $C$ of $B$. Its parent form is $\alpha$, so the same form-based
print and path maps are used for $C$ and $D$. The predecessor of the print
on $P_j(d)$ is an ordinary side of $B$, by the routing certificate. Its
third cell equals $C_j(d)$ by the child clause and equals $D_j(d)$ by the
first assertion. The print clause therefore supplies $D_j(d)$ for each
$j\le4$. This argument uses the predecessor in the small owner, rather
than an unsupported claim that the central tile is already legal.
\end{proof}

\begin{lemma}[Exact repaired decomposition]
\label{lem:comp/exact}
\lean{Tilings.Planned.macroblock_is_decomposition,Tilings.Planned.Strict.strict_macroblock_is_decomposition}
\uses{lem:comp/reconstruction,lem:comp/signals,def:control/substitution}
The finite patch $M$ is an exact member of $\widehat\sigma D$.
\end{lemma}
\begin{proof}
Check the clauses cell by cell. Child types and fixed first, second, and
fifth cells agree because each noncentral child has a parent of form $\alpha$.
Remaining inner nonprint fourth cells are $\bot$. Boundary fourth cells carry
$h=H(D)$. For an outer child side $b\subset\sigma'd$, its legal owner $B$
has a parent $C$ with $H(C)=h=H(D)$. Hence its choice between zero and its
fixed positive rank is exactly the choice prescribed by $D$. This is the
case missing from the unchanged composition argument in the source article:
the outer fourth value depends on the parent's mode, even though it is
not an ordinary copied cell.

Lemma~\ref{lem:comp/signals} verifies all ordinary third signals and all
four print cells in the noncentral owners. Remaining third cells are zero
by the child clause. Thus every noncentral tile has the required colours
as a child of $D$. (The outer-side owners used above have boundary sides
by Lemma~\ref{lem:comp/outer-boundary}.) Every side of the central tile is
a complete shared inner side with one of these legal noncentral neighbours. The neighbour forces
equality of the full colour by properness. The repaired substitution
prescribes equal colours to both owners of such a side, including the
borrowed tuple on a print. Consequently the central tile also has all the
required child cells. No initial legality assumption about the centre
has been used.
\end{proof}

\begin{lemma}[Cyclic reconstructed parents are legal]
\label{lem:comp/cyclic-legal}
\lean{Tilings.Planned.reconstructed_cyclic_legal,Tilings.Planned.Strict.strict_reconstructed_cyclic_legal}
\uses{lem:comp/exact,def:certificate,def:control/legal,lem:finite-idempotent}
Suppose the centre $A$ of $M$ is legal and the type $t$ of $D$ is cyclic:
it is the central child type in the template of its own form. Then $D=A$
as typed decorations, and $D$ is legal.
\end{lemma}
\begin{proof}
Both tiles have type $t$. The retained large form is fixed by the central
form map, obtained using Lemma~\ref{lem:finite-idempotent} in the geometric
construction. Its print map is the identity on marked sides. Hence every
marked first-four cell of $A$, as central child of $D$, copies the
corresponding cell of $D$. The remaining cells are local and agree because
both tiles are normal. The header is computed from the type and fourth
vector and therefore agrees as well. Thus $A=D$, proving the claim.
\end{proof}

\begin{lemma}[Noncentral reconstructed parents are legal]
\label{lem:comp/noncentral-legal}
\lean{Tilings.Planned.reconstructed_noncentral_legal,Tilings.Planned.Strict.strict_reconstructed_noncentral_legal}
\uses{lem:comp/exact,lem:comp/cyclic-legal,def:certificate,def:control/legal}
If the centre $A$ is legal and $t$ is noncentral, then $D$ is legal.
\end{lemma}
\begin{proof}
A normal noncentral tile has a marked ordinary inner side with a nonzero
first identity cell, by the separating certificate. Its print in $A$
contains that nonzero value. A legality chain of $A$ with only central
types would copy this marked first cell forever: all parent forms would
be retained large forms, whose print maps are identities. The cell would
have no origin and would be zero, a contradiction. Therefore the chain
reaches a legal noncentral ancestor.

Choose a shortest central-child chain ending at such an ancestor. Passing
to the central child fixes large forms. All but possibly the last two
vertices are therefore cyclic tiles of the same type and, by the argument
of Lemma~\ref{lem:comp/cyclic-legal}, identical decorations. Delete these
repetitions. The chain has length $n\le2$. Length two is impossible here:
the immediate parent of $A$ would have a large form, whereas its known
parent $D$ has a noncentral form. A child's position determines its parent
form, so these cannot both be parents. Thus $A$ has a legal noncentral
parent $E$ immediately above it.

The parents $E,D$ have the same form. Their print maps depend only on this
form, so their marked first-four fields agree through the common child $A$.
Let $\Phi$ be the multiset of type components of first-cell interface identitys
on marked sides. We have $\Phi(E)=\Phi(D)$. The noncentral fingerprint
certificate implies that their types agree. Their local fields then agree
by normality, while their marked first-four fields already agree. Hence
$E=D$, and $D$ is legal.

The fingerprint condition can be checked from exact finite data: original
types have a nonempty inner identity set; a template's connected dual graph
has more than two vertices; small types have enough ordinary marked inner
sides to separate their cardinalities from original types; and each pair
of small types has at least three private ordinary original-neighbour
occurrences. The exceptional special side contributes at most two arbitrary
components. These are finite separating conditions, not informal estimates
that sufficiently large quasi-squares will automatically have them.
\end{proof}

\begin{lemma}[Positive ranks align an entire macroside]
\label{lem:comp/alignment}
\lean{Tilings.Planned.alignment_complete}
\uses{def:certificate,def:tiling,def:control/substitution}
Let the macroblocks of a proper plane tiling of legal typed tiles form a
plane partition with disjoint interiors (legality makes every outer fourth
cell zero or the generated rank of its side; without this a constant
colouring satisfies the hypotheses below on the three-block counterexample).
If every child side of one macroside has a complete opposing side with the
same positive alignment rank, \emph{and every macroside of an opposing
macroblock is mode-uniform} (either all of its child sides carry positive
ranks, or all carry zero), then this macroside has a complete opposing
macroside with the same ordered subdivision.

Without the mode-uniformity hypothesis the statement is false: let $d$
have children of ranks $1,2,3$, each completely matched by the children of
ranks $1,2,3$ of an opposing macroside $e$ which has five children, the
last two facing trivial children of a third block. The hypothesis on $d$
holds, yet $d$ has no complete opposing macroside. This example is a
counterexample only under the weaker legality ``every outer fourth cell is
zero or the generated rank of its side''; under the full-header legality of
the repaired child relation every macroblock of a proper legal tiling has a
common boundary header, so mode-uniformity follows from legality and the
hypothesis is automatic. In
Lemma~\ref{lem:comp/proper} mode-uniformity is available, because every
block is an exact decomposition of its reconstructed parent and the parent
header fixes one mode per side.
\end{lemma}
\begin{proof}
Use the global orientation of ranks along straight sides, required by (G3).
The child side
of rank one coincides completely with a child side of rank one on an
opposing macroblock. Their first endpoints, and thus the beginnings of
their macrosides, coincide. Continue along the finite ordered list of child
segments. An opposing macroblock cannot end inside one of these segments:
the segment has a complete matching side. If it ends between consecutive
segments, the next opposing macroside would have to begin there and its
first rank would be one, whereas the next given rank is larger than one.

The other end needs the mode-uniformity hypothesis. Let the given side
have $L$ children and the opposing side $e$ have $L'$ children, and suppose
$L'>L$. Since child $1$ of $e$ carries a positive rank, mode-uniformity
makes child $L+1$ of $e$ ranked, hence nontrivial, so it needs a complete
match of rank $L+1$. Beyond the end of the given side, on the same side of
the line, lies either the next listed side of the same block, whose first
child has rank $1$, or a side of a different block, which would have to
begin $L$ children earlier and overlap the given block, or a trivial side,
which cannot match a nontrivial one. All three are impossible, so $L'=L$.

For completeness, an alternative opposing macroside cannot have begun
earlier and merely join the contact at this switch. An open portion of that
earlier side would lie opposite the same open portion of the original
opposing block, with their interiors on the same half-plane side. Local
polygonal interior neighbourhoods would overlap, contradicting the
disjointness of macroblock interiors. Coverage supplies an opposing block
after each endpoint, so the rank-reset obstruction cannot be avoided by
leaving a gap. Applying the argument at the other end shows that the
opposing side terminates exactly with the given side. Complete child-side
matching then identifies the whole ordered subdivision. This geometric
lemma must be formalized separately from the algebraic header kernel.
\end{proof}

\begin{lemma}[The reconstructed plane tiling is proper]
\label{lem:comp/proper}
\lean{Tilings.Planned.composed_proper,Tilings.Planned.Strict.strict_composed_proper}
\uses{lem:comp/assembly,lem:comp/exact,lem:comp/alignment,lem:comp/reconstruction,lem:comp/boundary-header}
For a proper legal plane tiling $\mathcal T$, replace each assembled
macroblock by its reconstructed parent and scale down by the substitution
factor. Each assembled macroblock is replaced by its reconstructed parent in
the sense of Lemma~\ref{lem:comp/reconstruction}, i.e.\ the parent whose
type is the type of the common boundary header of the block given by
Lemma~\ref{lem:comp/boundary-header}. The result $\mathcal T'$ is a proper tiling of normal identity-correct decorations
with active inner sides, and $\mathcal T$ is an exact repaired decomposition
of $\mathcal T'$.
\end{lemma}
\begin{proof}
Assembly partitions the plane into the macroblocks. Their underlying
polygons are scaled parent polygons, so replacing and scaling them preserves
coverage and disjoint interiors. Exact decomposition follows block by block
from Lemma~\ref{lem:comp/exact}; normality, identity-correctness, and activity of inner sides follow
from Lemma~\ref{lem:comp/reconstruction}.

Every macroside of every block is mode-uniform: the block is an exact
decomposition of its reconstructed parent $E$ (Lemma~\ref{lem:comp/exact}),
and the outer fourth cells on $\sigma'e$ are all zero or all positive
ranks according to $z_E(e)$. So the alignment lemma applies.

Consider a nontrivial side $d$ of a reconstructed parent $D$. If it is
inner, $z_D(d)=0$ by definition. If it is outer, its fifth cell is zero,
so nontriviality gives $D_4(d)\ne0$ and again $z_D(d)=0$. Thus all its
child outer fourth cells are positive ranks. Properness of $\mathcal T$
matches each completely with an equal positive rank. The alignment lemma
gives a complete opposing parent side and the same subdivision. Its five
ports consequently coincide with the corresponding five ports on the
given side. The third cells on those ports agree in $\mathcal T$. These
five values are exactly the five reconstructed parent cells, proving full
colour equality. A trivial outer side imposes no requirement. Hence
$\mathcal T'$ is proper.
\end{proof}

\begin{lemma}[Central acyclic reconstructed parents are legal]
\label{lem:comp/acyclic-legal}
\lean{Tilings.Planned.reconstructed_acyclic_legal,Tilings.Planned.Strict.strict_reconstructed_acyclic_legal}
\uses{lem:comp/proper,lem:comp/exact,lem:comp/noncentral-legal,lem:comp/cyclic-legal,def:certificate}
Every reconstructed parent of central acyclic type in $\mathcal T'$ is legal.
\end{lemma}
\begin{proof}
Let $D$ be such a parent, of type $t$, and let $A$ be its original legal
central child. Every noncentral or cyclic tile of $\mathcal T'$ is already
legal by the preceding legality lemmas. All reconstructed tiles are normal
and identity-correct with active inner sides, so the assembly lemma also applies to
$\mathcal T'$. In the second-level macroblock containing $D$, every
noncentral tile is legal. The full repaired reconstruction and exactness
lemmas give a normal parent $C$ of type $r$, with $D$ its central child.
No unproved inheritance of a single-scale flag rule is needed.

The type $r$ is noncentral. If it were central, its large form would be
fixed by the central-child form map. It and its central child $t$ would
have the same large form, and hence the same central child type, making
$t$ cyclic. This contradicts the assumption on $D$.

The nonzero marked first cell of normal noncentral $C$ is copied into
$D$ and then into $A$. The origin argument from the noncentral legality
lemma therefore shows that a legal ancestry chain of $A$ reaches a
noncentral type; removing cyclic repetitions gives length at most two.
An immediate parent of $A$ has the same large form as $D$, and hence a
central type. The chain must have length exactly two:
\[
 A\ \longleftarrow\ E\ \longleftarrow\ F,
\]
where $E,F$ are legal and $F$ is noncentral.

The parents $E,D$ have the same form and agree in their marked first-four
fields through $A$, so $\Phi(E)=\Phi(D)$. Write $L=4N$ for the number
of marked sides of a large form, and $Q$ for the maximum on a noncentral
form. At most $Q$ marked sides of $D$, the central child of $C$, are prints.
On each of the remaining at least $L-Q$ sides its first interface identity
contains $t$. Thus the multiplicity of $t$ in $\Phi(D)$ is at least $L-Q$.
Suppose $E$ had a different central type $t'$. Its nonprint identities
contain $t'$ and a small neighbour type, neither equal to $t$. Only its
at most $Q$ prints can contribute $t$; a copied print cell contributes at
most one occurrence of $t$ per print under the type projection, so the
multiplicity is at most $Q$, and in any case at most $2Q$.
This contradicts $L-Q>2Q$, supplied by the certificate. Therefore $E$
also has type $t$. Normal local cells and the marked copied cells now
identify $E=D$. Since $E$ is legal, $D$ is legal.
\end{proof}

\begin{theorem}[The repaired composition theorem]
\label{thm:compose}
\lean{Tilings.Planned.compose_legal_exists}
\uses{def:certificate,def:hereditary,lem:comp/proper,lem:comp/noncentral-legal,lem:comp/cyclic-legal,lem:comp/acyclic-legal}
For every macrotemplate system satisfying Definition~\ref{def:certificate}
with hereditary listed sides (Definition~\ref{def:hereditary})
and the explicit header substitution of Definition~\ref{def:control/substitution},
every proper legal plane tiling has a proper exact repaired composition
whose tiles are legal.
\end{theorem}
\begin{proof}
Construct $\mathcal T'$ by assembly and reconstruction. The properness
lemma proves that it is a plane tiling, is proper, and decomposes exactly
to the input. Its reconstructed types are either noncentral, central cyclic,
or central acyclic. The three legality lemmas cover these disjoint cases.
Thus every tile of $\mathcal T'$ is legal. The conclusion uses the same
properness and legality predicates as the input, so the theorem can be
iterated. The header modification is what provides this closure.
\end{proof}

\begin{theorem}[Lifting an actual side-to-side-supertiling hierarchy]
\label{lem:lift-hierarchy}
\lean{Tilings.Planned.lift_hierarchy}
\uses{def:certificate,def:hereditary,def:hierarchy,def:control/origin,def:control/substitution}
Suppose the macrotemplate system has hereditary listed sides, so that no
fine outer side straddles a parent listed vertex and ``the unique
containing side'' below is defined. Suppose an actual $\sigma'$ hierarchy
is given and its finite supertiles
are side-to-side at every level: for every tile $X$ of level $n$ and every
$k\le n$ the patch $\sigma'^{\,k}X$ of level-$(n-k)$ tiles is side-to-side
(the full decomposition alone is not enough: two level-$1$ siblings with
partially overlapping sides would give coincident fine sides different port
ranks). Then it admits compatible legal decorations at every
level under the repaired substitution; its decorated level-zero tiling
is proper. This side-to-side-supertiling assumption is an additional
geometric instantiation obligation, supplied by a correctly grouped STS
construction, not proved by the header kernel.
\end{theorem}
\begin{proof}
At each level the hierarchy specifies each child's type and its parent.
On an outer side, follow the unique containing side through parent levels.
If it stays outer forever, put its fourth cell equal to zero at all levels
of that chain. If it eventually becomes inner, assign the fixed positive
alignment ranks on the outer levels below the first inner side. These
choices obey the outer-fourth substitution clause. Set remaining inner nonprint
fourth cells to $\bot$, and boundary fourth cells to the computed header
of their parent.

Set every remaining noncopied fresh cell to its prescribed child value;
this includes fixed local cells and first identities on special nonprints.
For a recursively copied cell, follow its source: use the origin value if an origin is reached, and
zero if it is not. Copy-source uniqueness makes this assignment
unambiguous. Equality along a copy edge holds in both cases. A header depends
only on the type and the outer fourth cells, all of which were assigned
before the copying step: inner entries of its mode vector are false regardless
of their fourth values. Its assignment therefore introduces no additional
recursive choice. The child clauses are now
all satisfied. Each actual ancestor chain witnesses legality, including
the no-origin condition at every level.

Inside a macroblock, shared side colours agree because both owners receive
the same parent data; synthesized boundary headers also agree. For an outer
side with nonzero fourth cell, its ancestor side becomes inner at a finite
level. Both overlapping sides therefore belong to one finite higher
supertile. The assumed STS property makes their descendant sides coincide
completely, with matching ranks and ports. Their copied cells either reach
the same shared origin, or matching origins on a shared side, or both have
no origin and are zero. The second alternative is an induction on the
level: the two copied cells on the coincident sides $b\subset\sigma'd$ and
$b'\subset\sigma'd'$ copy the cells of the parent sides $d,d'$, which are
again coincident sides in the next level with one fewer level to go before
the chain becomes inner; the induction hypothesis gives equality there.
The third cell of an outer side is nonzero only when that side is a port,
by the inner-side clause of (G4); hence no outer side other than a port
carries a parent signal, and the coincident side of a port is the port of
the same number of the opposing block. Hence their full colours agree. The same argument
applies to inner sides, already lying in a finite supertile. An outer side
whose chain stays outer forever has zero fourth and fifth cells and is
trivial, so its shifted contacts are permitted. The level-zero decoration
is proper, and the construction is compatible across all hierarchy levels.
\end{proof}

\begin{theorem}[Soundness and the conditional representation]
\label{thm:comp/soundness}
\lean{Tilings.Planned.projects_to_hierarchy}
\uses{thm:compose,lem:lift-hierarchy,def:hierarchy,def:certificate,def:hereditary}
Under the certificate assumptions, with hereditary listed sides, erasing
the colours of every proper
legal tiling gives a $\sigma'$-hierarchical tiling. If in addition the
actual hierarchies to be represented have side-to-side finite supertiles,
then the erasures of proper legal tilings are exactly those hierarchical
tilings. The resulting decorated tile alphabet is finite.
\end{theorem}
\begin{proof}
Apply the repaired composition theorem repeatedly. Classical dependent
choice, or recursive selection of a witness for each natural-number level,
gives proper legal tilings $\mathcal T^{(0)},\mathcal T^{(1)},\ldots$,
each an exact composition of its predecessor. Erasure commutes with the
repaired substitution, since the underlying geometry was unchanged.
Thus their erased sequence witnesses the required $\sigma'$ hierarchy.
This direction uses no assumed existence of a coloured hierarchy.

Conversely, the lifting theorem supplies a proper legal decoration of
every hierarchy satisfying its geometric STS condition. Finiteness follows
because there are finitely many forms, types, side positions, and ordinary
index tags, while the dependent header sum is finite. Legality is a subset
of this finite alphabet. These two implications prove the conditional
representation statement.

This theorem is deliberately conditional on the geometric certificate and
the stated realization condition. It is not a completed proof of the
general STS/FLC soficity theorems. Establishing a valid connected grouping,
the routing and separation certificates, the compatibility with the
original substitution, and the repaired FLC geometric reduction remains
outside this chapter. In the crown-bearing variants, positive outer
fourth cells make every header mode vector false; the header then carries
the former parent type with redundant control bits. Crown consistency and
the additional language argument require their own nodes and are not
inferred from the conditional hierarchical representation proved here.
\end{proof}

\section{From the repaired construction to the main theorems}

The grouping chapter now gives ordinary mathematical proofs of the
quantitative localization and uniform routing lemmas. The polygonal Jordan
and PL annulus facts used there remain substantial Lean library work.
No main theorem below has been checked by Lean. Compactness is applied only
to geometric group/clone labels. After obtaining an actual infinite
hierarchy, compute zero statuses, headers and copy-origin values in that
order. Origin-based legality is not asserted to be closed under arbitrary
limits of finite coloured hierarchies. We distinguish the
hereditary-side core from the original hypotheses; the FLC reduction and
endpoint schedules bridge that distinction without assuming the result.

\begin{lemma}[Hierarchy with allowed crowns is substitutive]
\label{lem:allowed-hierarchy}
\lean{Tilings.allowed_crowns_hierarchy_is_substitutive}
\uses{def:hierarchy,def:flc,lem:crown-cover-plane}
Suppose $\tau$ has FLC and $T=T_0,T_m,T_{2m},\ldots$ is a $\tau^m$-hierarchy
whose crowns at every listed level are allowed. Then $T$ is a substitution
tiling.
\end{lemma}
\begin{proof}
Given a finite patch $F\subset T$, choose $k$ with
$\theta^{-km}\operatorname{diam}|F|<\varepsilon$.
The inverse-scaled support is covered by a crown $C$ of $T_{km}$.
All the level-zero tiles of $F$ belong to descendants of this crown:
coverage by the descendants and disjoint interiors give actual tile
inclusion, not only support inclusion, by
Lemma~\ref{lem:covered-tile-belongs}. The allowed crown $C$ occurs in
some supertile $S$. Thus $F\subset\tau^{km}S$, which is a supertile.
\end{proof}

\begin{lemma}[Crown colours enforce actual allowed crowns]
\label{lem:crown-colours}
\lean{Tilings.Decoration.crown_indices_sound}
\uses{def:certificate,def:crown,def:flc}
In the hereditary-side crown-coloured core construction, the uncoloured scattered
tiling at every reconstructed level has only allowed crowns.
\end{lemma}
\begin{proof}
Inside a macroblock the scattered tiles are a patch of an original
supertile, so full crowns away from its boundary are allowed.
Only generated fields at macrotemplate-boundary vertices are used here.
A copied print-second field is not asserted to be a crown at its physical
endpoint; prints are strictly inside the template.
At a macroboundary vertex every incident border tile records the crown
at that vertex. Legality gives the within-tile consistency of records on
two sides with that vertex; side matching gives equality across tiles.
The finite cyclic incident fan is connected through these sides, so all
records identify one allowed crown. Each incident original tile is included
in that record, including every constituent of a quasi-tile at the vertex.
The ambient fan is complete. An additional recorded tile would have an
open angular sector overlapping this complete fan, contrary to disjoint
interiors. Hence the recorded crown is exactly the ambient full crown.
This argument needs full pointed crown records, not merely a list of
neighbour colours or unpointed polygon shapes.
\end{proof}

\begin{lemma}[Dissection and glue colours]\label{lem:glues}
\lean{Tilings.Decoration.dissection_preserves_sofic,Tilings.Planned.Strict.exists_coincident_side}
\uses{def:certificate,def:localrule}
If the added quasi-forms are finite connected polygonal unions of original
tiles with connected dual graphs, their coloured local-rule presentations
can be transferred back to a finite decoration of the original forms.
\end{lemma}
\begin{proof}
For every added coloured form and every internal constituent side,
introduce a unique glue carrying the form identity and the two constituent
positions. Keep the old colour on its exposed boundary segments.
In a proper scattered tiling, each glue forces the rigidly positioned
partner. Connectedness of the finite constituent dual graph forces the
whole added tile; uniqueness of position tags prevents competing copies.
Conversely, scattering a proper quasi-tile tiling gives these matching
glues. Thus assembly and scattering are inverse on the presentations.
All checks involve tiles within a uniformly bounded diameter, since there
are only finitely many added forms. After every final transfer restore
surjectivity onto any unused original forms by inert decorations. In the
generic presentation forbid their singleton tiles. In the strict
presentation give an owner-side incidence a fresh colour whose only possible
matching incidence would force coincident interiors. These decorations
occur in no proper plane tiling and change no projected family.
The disconnected assignment counterexample fails exactly at this
connected-dual hypothesis.
\end{proof}

\begin{definition}[Finite endpoint schedules]\label{def:endpoint-schedule}
\lean{Tilings.Alignment.EndpointSchedule}
\uses{def:certificate,def:substitution,def:sts}
Fix the simulation power $m$, $s=\theta^m$, and one positive unit direction
for each unoriented supporting-line direction. Every whole fine outer side
$e$ has coordinates $[0,|e|]$ from its first endpoint in this direction.
For a parent group $A$ use the boundary layers of $sA$ with its listed sides,
and of $\theta^{m-j}\tau^j\Delta A$, $j=0,\ldots,m-1$.
In each layer, intersect its listed boundary sides with $e$. Each positive
atom $[u,v]$ has a unique containing coarse listed side with coordinates
$[p,q]$; record $(u,v;p,q)$, allowing $p,q$ outside $e$.
The layered ordered tuple is $W(A,e)$. Its tag is the geometric data,
not the owner form or edge ID: different forms with equal data share a tag.
The layer index ($G$, or $j=0,\ldots,m-1$) is itself a uniform tag that
does not depend on the form, so that schedules of different parent forms
are compared layer by layer. Layer $G$ records the same data as layer
$j=0$; it is kept as an explicit safeguard and is harmless.
The alphabet is finite because the templates, $m$, and their sides are finite.
No fine side is physically split and no repeated scaling orbit is stored.
\end{definition}

\begin{lemma}[Schedules align all intermediate levels]\label{lem:alignment-tuple}
\lean{Tilings.Alignment.schedule_intermediate_sts}
\uses{def:endpoint-schedule,def:sts}
After a proper fine STS tiling is assembled into the fixed macrotemplates,
equal schedules on shared whole outer fine sides force the grouped parents,
their scattered original tiles, and every inserted original level to be STS.
\end{lemma}
\begin{proof}
Choose any positively overlapping coarse listed sides $c,d$ on opposite
macroblock boundaries, in the same layer tag. Pick a point of their overlap
away from the finite fine and schedule endpoints. The whole fine sides
through it coincide by fine STS. Their canonical origins and directions
therefore coincide. Equal schedules give the same atom through that point
and the same two coordinates $p,q$ for its coarse containing side.
Thus $c,d$ have exactly the same endpoints. Inside a block each original
layer is already an STS supertile patch. This proves every layer separately,
even at a persistent interface with no common higher ancestor.
A genuine corner cannot lie inside a fine side: its wedge must be covered
by child sectors, and STS excludes a partial overlapping edge at that point.
Only collinear dummy endpoints can be interior schedule cuts.
\end{proof}

\begin{lemma}[Whole-side ports without nested subdivisions]\label{lem:whole-ports}
\lean{Tilings.Alignment.whole_side_ports}
\uses{def:mesh,def:endpoint-schedule,lem:alignment-tuple}
For each parent listed side $a$, order the complete fine outer sides whose
closed segments lie strictly inside $sa$, in the same global direction.
This list has at least $(s|a|-2D)/D$ entries. Its five interior sixth-quantiles
$\lceil iL/6\rceil$, $i=1,\ldots,5$, give five distinct ports for large $s$. After schedule
alignment, corresponding parent sides have the same physical port lists.
\end{lemma}
\begin{proof}
Only the at most two fine edges meeting the endpoints of $sa$ are omitted;
their total covered length is at most $2D$. The count uses no side-to-side
hypothesis: the at most two fine sides covering pieces of the parent side
near its endpoints---an outer side across a dummy vertex, or a non-outer side
through a reflex corner---are absorbed in the $2D$ term. Every remaining edge is at most
$D$ long, proving the count. The positive minimum fine-edge length gives
the quantile separation estimates with a fixed $O(D)$ error. All five ports,
including the singleton fifth one, stay uniformly separated from the listed
endpoints and from one another in normalized coordinates.
Once two parent sides coincide, fine STS gives identical whole-edge
subdivision of their common support; the endpoint exclusion is identical.
Number from the global supporting-line orientation, not opposite boundary
traversal orders. All five paths of an unmarked parent side remain singleton
ports; all five parent components are still transported.
\end{proof}

\begin{lemma}[Strict composition with schedules]\label{lem:schedule-compose}
\lean{Tilings.Alignment.strict_schedule_composition,Tilings.Planned.Strict.strict_opposing_macroside,Tilings.Planned.Strict.strict_compose_exists}
\uses{lem:whole-ports,lem:alignment-tuple,lem:comp/assembly,lem:comp/boundary-header,lem:comp/signals,lem:comp/cyclic-legal,lem:comp/noncentral-legal,lem:comp/acyclic-legal}
In the strict variant replace generated outer fourth ranks by tagged $W$.
Every proper legal tiling has an exact proper legal coloured composition,
and its scattered intermediate levels are STS.
\end{lemma}
\begin{proof}
Use disjoint tags for schedules, headers, identities, zero and $\bot$.
In this variant the outer fourth cell $W(A,e)$ of a legal fine tile is
fixed by its type, since legality makes the tile a child of its parent's
form; so it is a local cell here, unlike the flexible variant. The parent's
own fourth value is transported separately through port $4$; equality of
these transported values across a persistent fault line follows from
completeness one level up, with originless copied cells set to zero.
A schedule fourth cell with fifth zero can match only another outer side.
The assembly and fingerprint arguments use inner identities and are unchanged.
The boundary belt gives a common parent type/header; each parent's cells
are recovered from its five port third cells. Legal boundary witnesses
identify the template, so its fixed generated $W(A,e)$ agrees with the
reconstructed parent's decomposition independently of that parent's fourth
value. The other cells and the centre follow by the same path/print recovery.
Schedule alignment gives complete parent sides, and the whole-port lemma
then pairs equal numbered ports. Equality of their third cells identifies
all five parent components, including the transported fourth value.
Legality uses the same finite copy-origin argument and the cyclic,
noncentral and second-scale acyclic fingerprint cases. Schedule tags are
generated origins, replacing integer ranks without altering these cases.
This repeats the finite control proof with a different generated field;
it does not claim the two decorated substitutions are the same function.
\end{proof}

\begin{lemma}[Projection of a strict composition]\label{lem:strict-projection}
\lean{Tilings.Alignment.StrictScheduleSystem.isComposition_project}
\uses{lem:schedule-compose}
In a strict schedule system the projection of a coloured composition is a
grouped composition.
\end{lemma}
\begin{proof}
The coloured rules project onto the grouped rules and the expansions agree,
so the projection of the decomposition of $T''$ is the decomposition of the
projection of $T''$.
\end{proof}

\begin{lemma}[Strict presentation of an STS substitution]\label{lem:strict-presentation}
\lean{Tilings.Alignment.exists_strict_schedule_system,Tilings.Alignment.StrictScheduleSystem.HasCompositions,Tilings.Alignment.StrictScheduleSystem.Complete,Tilings.Planned.Strict.strictSystem,Tilings.Planned.Strict.strictSystem_hasCompositions,Tilings.Planned.Strict.strictSystem_complete,Tilings.Planned.Strict.strict_actual_schedules_aligned,Tilings.Planned.Strict.strict_lift_decorations,Tilings.Planned.Strict.strict_lift_matched,Tilings.Planned.Strict.exists_strict_schedule_system_of_certificate}
\uses{def:sts,lem:certificate-exists,lem:schedule-compose,lem:whole-ports,def:endpoint-schedule}
For every STS substitution $\tau$ there is a strict schedule system (the
strict typed decoration of the certificate of
Lemma~\ref{lem:certificate-exists}, with second component zero, endpoint
schedules and whole-side ports) whose grouped family contains every original
form as a singleton group and whose dissections remain side-to-side under
all decompositions, such that every proper strict coloured plane tiling has
an exact proper coloured composition, and every grouped $\sigma'$-hierarchy
all of whose levels and scattered intermediate original levels are
side-to-side is the projection of a proper strict coloured plane tiling.
\end{lemma}
\begin{proof}
Construct the finite groups on the original STS cellulations, retaining all
original listed vertices (Lemma~\ref{lem:certificate-exists}); every original
form is kept as a singleton group. Use strict matching, second component
zero, endpoint schedules and whole-side ports. Exact proper compositions are
the existence half of Lemma~\ref{lem:schedule-compose}. For completeness, at
each shared fine side $e$ between blocks $A$ and $A'$, a layer-$j$ side $c$
of $A$ on the boundary of $sA$ lies wholly on the interface with a single
neighbour: by STS of the actual level its far side is one layer-$j$ tile,
which lies in one block, namely $A'$. Hence the actual original intermediate
levels have identical endpoints on both sides of $e$, so their schedules and
port signals match. Origin chains give legal tiles.
\end{proof}

\begin{lemma}[Soundness of the strict presentation]\label{lem:strict-sound}
\lean{Tilings.Alignment.sts_hierarchical_of_strict_proper}
\uses{lem:strict-presentation,lem:strict-projection,lem:schedule-compose,lem:alignment-tuple,def:hierarchy}
For a strict schedule system with exact proper compositions whose
dissections remain side-to-side under all decompositions, the scattering of
the projection of a proper strict coloured plane tiling is a side-to-side
$\tau$-hierarchical plane tiling.
\end{lemma}
\begin{proof}
Repeatedly apply strict schedule composition. At each step insert the
original levels $j=0,\ldots,m-1$, all STS by schedule alignment. This yields
an infinite original $\tau$ STS hierarchy.
\end{proof}

\begin{lemma}[Completeness of the strict presentation]\label{lem:strict-complete}
\lean{Tilings.Alignment.exists_strict_proper_of_sts_hierarchical}
\uses{lem:strict-presentation,lem:grouping-lift,def:hierarchy}
For a complete strict schedule system whose grouped family contains every
original form as a singleton group, every side-to-side $\tau$-hierarchical
plane tiling is the scattering of the projection of a proper strict coloured
plane tiling.
\end{lemma}
\begin{proof}
Lift the given STS hierarchy by finite-height grouping compactness
(Lemma~\ref{lem:grouping-lift}, STS levels); all its levels and scattered
intermediate levels are the given STS levels. Completeness of the system
gives the proper strict coloured tiling.
\end{proof}

\begin{lemma}[Finite collar conversion]\label{lem:collar-conversion}
\lean{Tilings.LocalRule.side_presentation_of_sts_sofic}
\uses{def:localrule,def:sofic,def:sts,def:mesh,lem:localfinite}
A sofic family consisting entirely of original STS plane tilings is
side-to-side sofic on those original listed sides.
\end{lemma}
\begin{proof}
All tilings here cover the plane; this is used when the reverse direction
needs every side to have an opposite neighbour, and when inert decorations
are excluded. Let $r$ bound the old forbidden diameters, $D$ bound tile
diameters, and $a>0$ be the minimum area. A generic segment joining interiors of two
tiles in a bounded geometric patch crosses whole-side neighbours and avoids
vertices. All crossed tiles lie in a fixed ball of radius $r+4D$.
Disjoint interiors and area at least $a$ bound their number by a uniform
integer $K$ (increase it once for all endpoint perturbations); erase loops.
Thus any old forbidden patch meeting a chosen root tile fits in its
dual-graph ball of radius $K$.

There are finitely many embedded labelled graph balls of radius $K+1$:
side degree and labels are finite, and a whole-side match has finitely many
relative placements. A candidate collar is a connected STS patch rooted at
its old decorated tile, with every tile at graph distance at most $K+1$,
and every side of a tile at distance at most $K$ having its opposite
neighbour inside the patch. Retain only candidates containing no old
forbidden patch. Global extendibility is not required.
On a shared side $A,B$, its colour records both canonically side-anchored
embedded labelled $K$-balls, in the common left/right owner order. Each
ball is contained in either owner's $(K+1)$-collar. Equality identifies
old labels and placements on both restrictions.

Prove simultaneously for all actual roots and $k=0,\ldots,K+1$ that the
state's intrinsic $k$-ball is its actual old-labelled dual $k$-ball.
The case zero is the root label. Edge equality identifies each predicted
neighbour with the actual one, and identifies its $K$-ball in the two
states. For $k\leq K$, restrict this to its $k$-ball and use the induction
hypothesis. The identity
$B_{k+1}(A)=\{A\}\cup\bigcup_{B\sim A}B_k(B)$, valid inside a complete
collar at these depths, proves the next case including placements.
Thus every actual bounded old forbidden patch lies in a locally admissible
stored collar, contradiction. Conversely actual proper old tilings give
these consistent collars. Forgetting the old decoration yields precisely
the given family. For a form absent from all proper tilings add an inert
decoration with a fresh distinct colour on each side. Its only possible
same-colour mate is the same owner-side incidence of an identical translate,
forcing coincident interiors; it cannot occur in a proper plane tiling.
This supplies surjectivity onto otherwise unused original forms.
\end{proof}

\begin{theorem}[Theorem 1(a), STS substitution tilings]\label{thm:sts-substitution}
\lean{Tilings.Main.sts_substitution_sofic}
\uses{def:sofic,def:sts,lem:sts-flc,thm:flc-substitution,lem:collar-conversion}
If $\tau$ satisfies STS, its substitution plane tilings are side-to-side sofic.
\end{theorem}
\begin{proof}
STS gives FLC, so Theorem~\ref{thm:flc-substitution} gives a generic sofic
presentation. Every projected substitution tiling is original STS: each
finite pair occurs in an STS supertile. Apply finite collar conversion.
This proof retains original dummy vertices and avoids using a crown field
on a fine edge which straddles an ancestor's dummy endpoint.
\end{proof}

\begin{theorem}[Theorem 1(b), STS hierarchical tilings]\label{thm:sts-stshierarchy}
\lean{Tilings.Main.sts_side_to_side_hierarchy_sofic}
\uses{def:sofic,def:sts,lem:strict-presentation,lem:strict-sound,lem:strict-complete,lem:glues}
If $\tau$ satisfies STS, its original side-to-side hierarchical plane
tilings are side-to-side sofic.
\end{theorem}
\begin{proof}
By Lemma~\ref{lem:strict-presentation} take a strict schedule system with
exact proper compositions and completeness whose grouped family contains
every original form as a singleton group. The family of projections of its
proper strict coloured plane tilings is side-to-side sofic by definition,
and glue dissection (Lemma~\ref{lem:glues}) transfers this presentation to
the scatterings, on the original listed tile sides. By
Lemma~\ref{lem:strict-sound} every such scattering is a side-to-side
hierarchical plane tiling, and by Lemma~\ref{lem:strict-complete} every
side-to-side hierarchical plane tiling arises this way; hence the scattered
family is exactly the family of original side-to-side hierarchical plane
tilings.
\end{proof}

\begin{theorem}[Theorem 1(c), unrestricted hierarchical tilings]\label{thm:sts-hierarchy}
\lean{Tilings.Main.sts_hierarchy_sofic}
\uses{def:sofic,def:sts,lem:sts-flc,thm:flc-hierarchy}
If $\tau$ satisfies STS, its hierarchical plane tilings are sofic.
\end{theorem}
\begin{proof}
STS implies FLC; apply Theorem~\ref{thm:flc-hierarchy}. The hierarchy itself
may have shifted persistent interfaces and need not be STS.
\end{proof}

\begin{lemma}[Compatible grouped hierarchy lift]\label{lem:grouping-lift}
\lean{Tilings.Grouping.hierarchy_lift_iff}
\uses{def:certificate,def:hierarchy,lem:localfinite}
Original $\tau^m$ hierarchies lift to grouped $\sigma'$ hierarchies.
When the given original levels are STS, the lifted grouped levels are STS.
\end{lemma}
\begin{proof}
For a finite height start with original-form tiles at its top and apply the
fixed grouped substitution downwards, identifying scattered tiles by
$\Delta\sigma'=\tau^m\Delta$. Do not retain a separate ring independently
in every constituent: the grouped-parent rule retains only its chosen centre.
Each fixed original tile position has finitely many possible group labels
and positions. Diagonal finite-alphabet compactness over all countably many
positions and levels preserves every finite parent-child constraint.
The limit is compatible at all levels. Retaining every original boundary
vertex ensures original whole-side contacts induce whole grouped contacts.
Projection/scattering proves the converse immediately.
\end{proof}

\begin{theorem}[Hereditary STS core, substitution family]\label{thm:hereditary-substitution}
\lean{Tilings.Core.hereditary_sts_substitution_sofic}
\uses{def:hereditary,def:sts,lem:sts-flc,lem:certificate-exists,lem:substitution-hierarchy,lem:grouping-lift,lem:crown-colours,lem:allowed-hierarchy,lem:glues,thm:compose}
An STS substitution with hereditary listed sides has a side-to-side sofic
substitution family.
\end{theorem}
\begin{proof}
Use the repaired geometry, exact interface IDs, strict matching, crown
records and positive exterior ranks.  The crown-bearing rule is the strict
positive-rank rule of the composition chapter with one additional finite
component: at every macrotemplate-boundary vertex each incident border
tile records in its second cell a fresh crown record (the finite pointed
crown of Lemma~\ref{lem:crown-colours}); these generated records are
distinguished from copied print-second cells, are generated only at
boundary vertices, and are transported along port two like the other
finite data.  Legality is origin-based as in the strict rule, with the
crown records as additional generated origins.  Fix a certificate for a
power $\tau^m$ (Lemma~\ref{lem:certificate-exists}) with grouped
substitution $\sigma'$, expansion $\theta$ and compatibility
$\Delta\sigma'=\tau^m\Delta$; a child type $t$ has an owner form
$\operatorname{owner}(t)$, a form $\operatorname{form}(t)$ and an anchor
$a_t$ in the template $\sigma'\operatorname{owner}(t)$.  For a grouped
form $q$ let $V(q)$ be the set of original marked vertices of the scattered
dissection $\Delta q$.  Since $\tau$ is STS, it has FLC
(Lemma~\ref{lem:sts-flc}), and the allowed full pointed crowns whose
support covers a neighbourhood of their centre form a finite list up to
translation; we call them \emph{full crowns}.  The neighbourhood condition
removes the incomplete crowns occurring on the boundary of a finite
supertile.

\emph{Step 1 (finitely many marked vertices).}  For every grouped form
$q$ the set $V(q)$ is finite, being a finite union of the vertex sets of the
finitely many original tiles of $\Delta q$.  This elementary statement
involves no control rule and no tiling; it makes every vector of crowns
indexed by $V(q)$ a finite datum.

\emph{Step 2 (the crown catalogue and the child operator).}  A \emph{crown
vector} for $q$ is a family $\gamma=(\gamma_v)_{v\in V(q)}$ of full crowns,
and its \emph{collar} is
$\operatorname{coll}(\gamma)=\Delta q\cup\bigcup_{v\in V(q)}(v+\gamma_v)$.
The vector is \emph{jointly allowed} if a translate of the entire collar
lies in a single supertile $\tau^n p$; this prevents independently chosen,
mutually incompatible crowns at different vertices.  The \emph{crown
catalogue} $\Gamma(q)$ is the finite set of crown vectors whose collar has
disjoint interiors, is side-to-side, has exactly the pointed crown
$v+\gamma_v$ at every $v\in V(q)$, and is jointly allowed.  For a child type
$t$ and $\gamma\in\Gamma(\operatorname{owner}t)$ decompose the whole collar
by $\tau^m$ and take, at every $v\in V(\operatorname{form}t)$, the physical
star of $\tau^m\operatorname{coll}(\gamma)$ at $a_t+v$, translated back to
the origin.  Each such raw child star is again a full crown (Step~2a): the
collar covers an open neighbourhood of the support of $\Delta q$, because
side-to-side matching and the exact stars at the endpoints of every
boundary side saturate it, so by compatibility the decomposed collar covers
a neighbourhood of the child group $a_t+\Delta(\operatorname{form}t)$; the
star is allowed because the decomposed collar lies in the supertile
$\tau^{n+m}p$.  The resulting vector $F_t(\gamma)$ satisfies all four
catalogue conditions with the single joint witness $\tau^{n+m}p$, since the
child collar is a subpatch of $\tau^m\operatorname{coll}(\gamma)$
(Step~2b).  Hence $F_t$ is a map
$\Gamma(\operatorname{owner}t)\to\Gamma(\operatorname{form}t)$.

\emph{Step 3 (the packet channel in cell two).}  Order the two endpoints of
every grouped side by a fixed translation-invariant lexicographic order of
the plane.  For a tile of type $t$ carrying $\gamma\in\Gamma(\operatorname
{form}t)$, the \emph{packet} of its side $i$ is the ordered pair of optional
full crowns at the two endpoints $v$ of $i$: the entry is $\gamma_v$ if
$a_t+v$ lies on the macroboundary $\partial(\theta\operatorname{owner}t)$
and is absent otherwise.  The doubly absent packet is encoded by the old
zero symbol, which is what ordinary internal nonprint sides and cyclic
central chains carry.  A tile $A$ \emph{decodes} $\gamma\in\Gamma(\operatorname
{form}A)$ if every nonprint second cell of $A$ is the packet of $\gamma$
and, at every $v\in V(\operatorname{form}A)$ with $a_A+v$ off the
macroboundary, $\gamma_v$ is the crown of $\tau^m\Delta(\operatorname
{owner}A)$ at $a_A+v$.  Print second cells are deliberately ignored: they
are transit channels, not physical crown records, and prints lie strictly
inside the template.  The decoder returns $\gamma$ when it is unique and
otherwise fails; failure is rejected by the incidence colours of Step~4.
If $A$ decodes $\gamma$ and $t$ is a child type of $\operatorname{form}A$,
the generated child of type $t$ writes the packets of $F_t(\gamma)$ into its
nonprint second cells, and it decodes exactly $F_t(\gamma)$: at the
macroboundary its packets are those of $F_t(\gamma)$ by construction, and
off the macroboundary the star of the decomposed collar is the star of
$\tau^m\Delta(\operatorname{owner}t)$ because the whole neighbourhood lies
in the parent support.

\emph{Step 4 (the finite encoding).}  Enlarge the finite cell alphabet to
the disjoint sum of the old strict symbols, the packets, and the
\emph{extended headers} $(t,z,\operatorname{dec}(A))$, consisting of the
old positioned type, the zero-mode bits computed from the fourth cells, and
the decoded catalogue value; old symbols remain a separate summand, so a
packet cannot be confused with an identity, a schedule, zero, $\bot$ or a
header.  Extended tiles keep the old positioned type and five cells per
listed side.  The generated child of $A$ of type $t$ copies along the old
strict copying graph, gives prints their interface identity in the fifth
cell, gives inner nonprint sides their identity in the first cell, writes
the packet of $F_t(\operatorname{dec}A)$ in cell two of nonprint sides, the
schedule rank on outer sides, the extended header of $A$ on boundary sides
and $\bot$ on the remaining inner sides in cell four, and zero elsewhere.
An extended tile is \emph{founded-legal} if it has an infinite ancestor
chain in which every level is the generated child of the next one, every
level decodes, and every copied cell with no finite origin in the
reflexive-transitive closure of the copy-source graph is zero.  This finite
data is compiled into a finite side decoration by giving every
founded-illegal or inert form a unique self-matching colour on each side;
the resulting presentation satisfies three exact contracts: in a proper
tiling every tile is a founded-legal extended tile, distinct tiles with
coincident typed sides carry equal five-cell arrays there, and conversely a
side-to-side tiling by founded-legal extended tiles with matching arrays is
proper.  These contracts exclude the constant-colour degeneracy of an
earlier attempt.  The projected family of proper tilings is the control
family; it is side-to-side sofic by definition.  This step asserts no
reconstruction or completeness statement.

\emph{Step 5 (one-level reconstruction).}  The reconstruction contract
between a control tiling $\mathcal T$ and a grouped tiling $\mathcal T'$
consists of a proper colouring of $\mathcal T$, a colouring of $\mathcal T'$
which is side-to-side, decodes at every tile and matches on coincident
typed sides, the statement that $\mathcal T'$ is a $\sigma'$ composition of
$\mathcal T$, and the clause that every coarse tile $D$ has, for each child
type $t$ of its form, a fine tile at $\theta x_D+a_t$ equal to the
generated child of $D$ of type $t$.  Founded legality of the coarse tiles is
intentionally not part of the contract.  Every control tiling admits such
a parent: assemble macroblocks from interface IDs
(Lemma~\ref{lem:comp/assembly}), read one extended header $(t,z,\gamma)$ on
the connected boundary belt (Lemma~\ref{lem:comp/boundary-header}), and
reconstruct the parent cells from the third cells of the five ports on each
parent side (Lemma~\ref{lem:comp/reconstruction}).  Heredity makes each
parent listed side a union of whole fine sides, and the zero-mode header is
everywhere false on outer sides.  The reconstructed parent $D$ has the
nonprint second cells of the legal parents chosen at its ports, all of
which decode the common value $\gamma$, so $D$ decodes $\gamma$; generation
of every child of the block then follows from Lemma~\ref{lem:comp/exact}
together with Step~3, since a child off the macroboundary carries only
absent packets and its decoded value is forced by the template.

\emph{Step 6 (legality of reconstructed parents).}  Legality of each
reconstructed parent is proved before reconstruction is iterated, so the
argument is not circular.  The three cases of Theorem~\ref{thm:compose}
apply verbatim to the extended rule.  If the type of $D$ is central cyclic,
then $D$ is its own generated central child, and the constant chain
decodes at every level (Lemma~\ref{lem:comp/cyclic-legal}).  If it is
noncentral, the fingerprint identification of
Lemma~\ref{lem:comp/noncentral-legal} identifies $D$ with the legal parent
of a central child.  If it is central acyclic, one passes to the first
noncentral ancestor (Lemma~\ref{lem:comp/acyclic-legal}).  The crown
component does not disturb these identifications: unmarked sides and
central nonprint sides carry the absent packet, which is literally the old
zero symbol in cell two, and the decoded value of $D$ is part of the
extended header recovered at its ports.  Hence the reconstructed parent
tiling again belongs to the control family.

\emph{Step 7 (iterated reconstruction).}  Starting from a control tiling,
Steps~5 and~6 give a parent in the control family together with the
reconstruction contract, and dependent choice iterates this to a grouped
hierarchy $\mathcal T=\mathcal T_0,\mathcal T_1,\ldots$ of control tilings.
The contract witnesses between adjacent levels are retained; they prevent
the levels from silently choosing unrelated proper colourings, so that the
colouring of level $n+1$ is the one whose generated children are the tiles
of level $n$.

\emph{Step 8 (all-level allowed crowns).}  At every level $n$, every vertex
$w$ of the scattered tiling $\Delta\mathcal T_n$ has an allowed full crown;
this includes vertices on interfaces between macroblocks.  Inside a
macroblock the scattered tiles form the patch $\tau^m\Delta(\operatorname
{form}D)$ of a decoded collar lying in a supertile, so full crowns away
from its boundary are allowed.  At a macroboundary vertex every incident
border tile records, by Step~3, the crown at that vertex in the packets of
both of its sides through $w$; decoding gives within-tile consistency of
these two records, matching of coincident typed sides gives equality across
tiles, and the finite cyclic incident fan is connected through these
sides.  Hence all records identify one full crown, which contains every
incident original tile; an additional recorded tile would overlap the
complete fan, so by disjoint interiors the recorded crown is exactly the
ambient crown (Lemma~\ref{lem:crown-colours}).  It is allowed because the
catalogue consists of full crowns.

\emph{Step 9 (soundness).}  One decomposition by the power $\tau^m$ is the
$m$-fold iterate of the decomposition by $\tau$.  Compatibility
$\Delta\sigma'=\tau^m\Delta$ transports each exact grouped composition
$\mathcal T_{n+1}\to\mathcal T_n$ of Step~7 to a $\tau^m$ composition of the
scattered tilings, so $\Delta\mathcal T_0,\Delta\mathcal T_1,\ldots$ is a
$\tau^m$ hierarchy.  By Step~8 its crowns are allowed at every listed level,
and only then Lemma~\ref{lem:allowed-hierarchy} shows that the scattered
bottom tiling $\Delta\mathcal T$ is a substitution tiling.

\emph{Step 10 (completeness).}  Let $T$ be a substitution tiling.
(10a)~Repeatedly compose it by Lemma~\ref{lem:substitution-hierarchy} to a
hierarchy of substitution tilings $T=T_0,T_m,T_{2m},\ldots$, and lift the
$m$-spaced levels through grouping (Lemma~\ref{lem:grouping-lift}) to an
STS grouped $\sigma'$ hierarchy $\mathcal U_0,\mathcal U_1,\ldots$ with
$\Delta\mathcal U_0=T$.  To every tile $A$ of $\mathcal U_n$ attach the
\emph{genuine} crown vector, whose entry at $v$ is the actual full crown of
$\Delta\mathcal U_n$ at $x_A+v$; its collar is a patch of
$\Delta\mathcal U_n=T_{nm}$, and it is jointly allowed because every finite
patch of a substitution tiling lies in a supertile.  This catalogue value is
fixed before any control cell is assigned.  (10b)~Colour every level as in
Lemma~\ref{lem:lift-hierarchy}, choosing the record of the actual crown:
zero statuses first, headers second, then every cell by its finite origin
along the actual ancestor chain in the fixed hierarchy, or zero if it has
none.  The decoder returns the preassigned genuine vector at every
ancestor, since decomposing a genuine collar yields the genuine collars of
the children, i.e.\ $F_t$ of a genuine vector is genuine.  This founded
lift is tied to the one hierarchy of (10a), so different sides cannot
choose unrelated legal chains or unrelated crown vectors.  (10c)~For two
coincident typed sides of distinct tiles of one level and a coordinate
$k\in\{1,\ldots,5\}$ for which either ancestor chain has a finite source,
the two values coincide.  The cases are: sibling prints, whose two owners
receive the same parent cells and the same interface identity; sibling
nonprints, prescribed by the same template; external packets and schedules,
where both packets read the genuine crowns of the same physical endpoints
in $\Delta\mathcal U_n$ and the positive ranks agree as in the strict
construction; and the aligned port of a copied third cell, whose chain
reaches the first level at which the interface becomes inner, where the
two ports of the same number copy one shared parent cell.  (10d)~If neither
chain has a finite source, both values are zero by the origin-or-zero
clause, so all five cells match.  (10e)~The third contract of Step~4 turns
founded legality and this matching into properness of the bottom
colouring, whose projection $\mathcal U_0$ lies in the control family and
scatters to $T$.

\emph{Step 11 (the exact family).}  Steps~9 and~10 are the two inclusions:
the scattered control family equals the family of substitution tilings of
$\tau$.  The control family is side-to-side sofic by Step~4.  Finally
scatter the finite groups using glue colours (Lemma~\ref{lem:glues}).  This
is the core result applied to hereditary clones below, not an assumption
about original dummy vertices.
\end{proof}

\begin{theorem}[Hereditary STS core, arbitrary hierarchy]\label{thm:hereditary-hierarchy}
\lean{Tilings.Core.hereditary_sts_hierarchy_sofic}
\uses{def:hereditary,def:sts,lem:certificate-exists,lem:grouping-lift,lem:lift-hierarchy,thm:compose,lem:hierarchy-step,lem:glues}
An STS substitution with hereditary listed sides has a sofic arbitrary
hierarchical family.
\end{theorem}
\begin{proof}
Use zero crown component, the repaired header and the nontrivial-side rule.
The hierarchy-lift proof assigns zero exactly to permanently outer ancestry
statuses first, computes headers second, then propagates copied origins.
Heredity makes each fine side belong to a unique parent listed side; no
whole edge can straddle inconsistent parent zero modes.
Completeness follows from the lifted hierarchy and shared-higher-supertile
matching for every interface which eventually becomes inner. Permanently
outer interfaces have trivial colours and may shift. Exact proper legal
composition iterates; scattering and inserting intermediate levels gives
an original hierarchy. Glue transfer is a bounded generic local rule.
\end{proof}

\begin{lemma}[Finitely many supertile neighbourhoods]\label{lem:finite-neighbourhoods}
\lean{Tilings.FLC.finite_neighbourhoods}
\uses{def:flc,lem:crown-cover-supertile,lem:finiteparents}
Under FLC there are finitely many tile neighbourhoods in supertiles up to
translation, where a neighbourhood contains every tile sharing any point
with the selected tile.
\end{lemma}
\begin{proof}
All neighbourhood supports have diameter at most $3D$, where $D$ bounds
the original tile diameters. Choose $k$ with
$3D\theta^{-k}<\varepsilon$. In a supertile $\tau^n p$ of order $n\geq k$,
the inverse-scaled support of the selected neighbourhood lies in the
support of $\tau^{n-k}p$ and is covered by one of its full pointed crowns,
by the corrected supertile version of the crown lemma. Decomposing this
crown $k$ times covers the entire neighbourhood, and
Lemma~\ref{lem:covered-tile-belongs} turns coverage into membership of
every neighbourhood tile in the decomposed crown. There are finitely many
allowed crowns and each has a finite $k$th image, with finitely many
selected subpatches and distinguished tiles. Orders $n<k$ have finitely
many possibilities directly. Consequently the neighbourhood catalogue is
finite.
This proof needs the supertile version, not only the plane version, and
includes incomplete crowns.
\end{proof}

\begin{definition}[Genuine-corner collared clones]\label{def:clones}
\lean{Tilings.FLC.GenuineCornerClone,Tilings.FLC.refinedSubstitution}
\uses{lem:finite-neighbourhoods,def:substitution}
First erase original collinear listed vertices, keeping each original form
label and geometric carrier; coarsened forms with identical carriers but
different original labels remain distinct forms. FLC survives because crowns at genuine vertices
are obtained from the old finite catalogue by forgetting listed cuts.
A clone records such a form $A$ and its full realized supertile neighbourhood
$S_A$, including incomplete boundary contexts and singleton order-zero ones.
Add to the boundary of $A$ only genuine nonstraight vertices of underlying
neighbour carriers in $S_A$. Do not add their artificial clone vertices.
A child $B$ receives context $(\tau S_A)_B$. The finite catalogue is taken
for the carrier-labelled coarsened system, not for arbitrary plane contacts.
\end{definition}

\begin{lemma}[Clone STS, well-definedness and heredity]\label{lem:clone-sts}
\lean{Tilings.FLC.refinedSubstitution_sts,Tilings.FLC.refinedSubstitution_hereditary}
\uses{def:clones,def:sts,def:hereditary,lem:macrocorner}
The clone substitution is finite, commutes with forgetting labels, has STS,
and has hereditary listed sides.
\end{lemma}
\begin{proof}
Touching children have touching original parents, so the full one-ring
context determines every child context. Here ``listed'' for a clone vertex
$v$ of $B$ means $v\in\partial B$, so $v$ lies on an edge of $B$; and the
witness child used below lies in the decomposed context of $B$ also when
the witness tile is $B$'s own parent $A$ or a context tile touching $A$
only at a point. It is realized in the decomposed
witnessing supertile. Iteration gives the actual neighbourhood in the
iterated ambient witness, including needed boundary neighbours.
On overlapping sides, endpoints and further cuts are genuine vertices of
underlying incident carriers visible in both contexts. Hence the resulting
subdivision agrees on both sides, proving clone-supertiles STS.

Every clone vertex $v$ has a witness genuine corner of a context tile $C$.
Its image $\theta v$ is a genuine corner of $\theta C$, hence a genuine
corner of some child of $C$ by the local halfplane argument. Every child of
$A$ incident to that point sees this witness child in its new context and
retains the cut. Thus every scaled clone listed vertex is a child clone
vertex; the sides are unions of whole child sides. This proof adds no
irrational orbit of old dummy endpoints.
\end{proof}

\begin{lemma}[Substitution tilings lift to clones]\label{lem:clone-substitution-lift}
\lean{Tilings.FLC.substitution_clone_lift_iff}
\uses{def:clones,lem:clone-sts,def:hierarchy}
A plane tiling is a $\tau$-substitution tiling exactly when it is the
projection of a refined-substitution tiling.
\end{lemma}
\begin{proof}
Projection of a finite refined supertile is an original supertile, proving
one direction. In the other, label each tile of the original plane tiling
by its actual full neighbourhood. Each is realized in a supertile:
take a finite larger patch containing the neighbourhood and its incident
context. Given any finite collared patch, surround it by a finite original
patch containing it strictly in the interior. This larger patch occurs in
some original supertile $S=\tau^n p$. The neighbourhoods of the selected
interior tiles agree with those in $S$. The singleton clone $(p,\{p\})$
exists, and its $n$th refined image contains the selected collared patch.
Thus the lifted plane tiling is substitutive.
\end{proof}

\begin{lemma}[Arbitrary hierarchies lift to clones]\label{lem:clone-hierarchy-lift}
\lean{Tilings.FLC.hierarchy_clone_lift_iff}
\uses{def:clones,lem:clone-sts,def:hierarchy,lem:localfinite}
A plane tiling is $\tau$-hierarchical exactly when it is the projection of
a refined-substitution hierarchy.
\end{lemma}
\begin{proof}
Projection gives one direction immediately. Fix an original hierarchy for
the other. For each finite height $n$, label its top-level tiles by
singleton clones and decompose the labels downwards for $n$ steps.
Every resulting label is a realized clone. The values on the countable
set of all tiles at all levels lie in finite alphabets. For each finite
set of levels and tile positions, the compatibility conditions are
satisfied by every sufficiently high finite-height construction.
Finite-alphabet product compactness, or a finite inverse-system argument,
therefore gives a simultaneous infinite assignment respecting every
decomposition. It need not record the entire ambient neighbourhood across
a permanent hierarchical interface; partial realized clone contexts are
allowed. No invalid assertion that every hierarchical crown is allowed is
used.
\end{proof}

\begin{lemma}[Forgetting extra vertices preserves sofic presentations]
\label{lem:segment-rules}
\lean{Tilings.FLC.forget_vertices_preserves_sofic}
\uses{def:localrule,def:sofic,def:clones}
A finite presentation on refined sides translates to a bounded local rule
on the original carriers by colouring their finitely many virtual segments.
\end{lemma}
\begin{proof}
The transfer is stated for an arbitrary bounded local rule on clone tiles,
not only for side-colour rules: the core's hierarchical presentation also
contains glue and scattering rules. Since clone tiles share carriers with
the original forms, any bounded forbidden patch on clone tiles is a
bounded forbidden patch on original tiles with the clone data read as
colours. Keep the clone choice as finite tile colour and its virtual straight
boundary intervals and colours. A virtual interval may cross an old
collinear dummy vertex. The generic rule reads it as one interval; restoring
old listed cuts does not impose new segment-to-segment restrictions.
For strict matching require segment-to-segment coincidence and colour
equality. For the flexible hierarchical rule require it only for
nontrivial segments. These conditions are expressed by forbidden
two-tile configurations with overlapping segments. Full plane coverage
ensures that a nontrivial side has an opposite neighbour; forbidding
partial or mismatched overlaps forces a whole matching segment.
All interacting tile pairs have bounded diameter, since the carrier forms
are finite. The forbidden set can be infinite as translations across
interfaces vary continuously. Reinstating virtual vertices is inverse to
forgetting them at the level of these presentations.
\end{proof}

\begin{theorem}[Theorem 2(a), FLC substitution tilings]\label{thm:flc-substitution}
\lean{Tilings.Main.flc_substitution_sofic}
\uses{def:flc,def:sofic}
Every FLC substitution has a sofic family of substitution plane tilings.
\end{theorem}
\begin{proof}
\uses{lem:finite-neighbourhoods,lem:clone-sts,lem:clone-substitution-lift,thm:hereditary-substitution,lem:segment-rules}
Build the finite genuine-corner clone substitution; it has hereditary STS.
Apply Theorem\ref{thm:hereditary-substitution}. Lift and project with Lemma
\ref{lem:clone-substitution-lift}, and transfer the rule by
Lemma~\ref{lem:segment-rules}. Composition of the finite projections is
a finite carrier-preserving decoration of the original protoset.
No forbidden pattern at an unbounded scale is introduced by this transfer.
\end{proof}

\begin{theorem}[Theorem 2(b), FLC hierarchical tilings]\label{thm:flc-hierarchy}
\lean{Tilings.Main.flc_hierarchy_sofic}
\uses{def:flc,def:sofic}
Every FLC substitution has a sofic family of hierarchical plane tilings.
\end{theorem}
\begin{proof}
\uses{lem:finite-neighbourhoods,lem:clone-sts,lem:clone-hierarchy-lift,thm:hereditary-hierarchy,lem:segment-rules}
Use the same hereditary clone substitution and Theorem~\ref{thm:hereditary-hierarchy}.
The hierarchy lift is the compactness lift, not the actual-neighbourhood
lift for substitution tilings. Transfer the flexible segment rule back to
the original carriers. The geometric construction is supplied by the grouping chapter.
Clone hierarchies may have shifted permanent interfaces; only their
supertiles are required to be STS.
\end{proof}

\section{Implementation boundaries and acceptance criteria}

Keep geometry and finite control separate. A practical module order is:
\texttt{Geometry/PolygonalDisk}, \texttt{PlacedTile}, \texttt{Patch},
\texttt{PlaneTiling}, \texttt{Substitution}, \texttt{Crowns},
\texttt{Geometry/RegularGrouping}, \texttt{Geometry/Routing},
\texttt{Control/Indices}, \texttt{Control/Origins},
\texttt{Composition/Reconstruction}, \texttt{Composition/Legality},
\texttt{LocalRules}, \texttt{FLC/Clones}, and \texttt{Main}.

An intermediate certificate may state exact support equalities, connected
finite incidence graphs, valid finite routes, and finite fingerprint
separation. It must not contain \emph{a legal parent exists},
\emph{proper composition exists}, or \emph{the family is sofic} as a field.
Those are the theorems proved from it. The final public statements must
construct the certificate from the polygon and substitution hypotheses.

The regression examples (the old mixed exterior-flag example, the
point-attached quasi-square assignment, the negative sine example, the FLC
endpoint-gap example, the irrational-mark square (STS without heredity), and
the three-block counterexample to the alignment lemma without mode
uniformity) are not Lean targets: they are counterexamples about the
article's original rules or about dropped hypotheses, not parts of the proof,
and are kept as prose in the relevant lemmas (for the three-block example,
under the full-header legality of Definition~\ref{def:control/legal} mode
uniformity follows from legality, see Lemma~\ref{lem:comp/alignment}).
The final project must compile without
\texttt{sorry}, \texttt{axiom}, \texttt{admit}, \texttt{native\_decide}, or
an unadvertised stronger polygon assumption. \texttt{Classical.choice}
and the usual kernel axioms remain the ordinary mathlib foundations.

The current archive contains a checked header kernel, not the final Lean
proof of these main goals. Its independent certificate and proof statuses
are recorded in the manifest.

\section{Measured comparison and estimated Lean size}

The comparison snapshot is
\href{https://github.com/AlexeyMilovanov/kolmogorov-complexity-lean/tree/42df7b843d6727291dfd0cbc1ceb299723ce5cb6}
{AlexeyMilovanov/kolmogorov-complexity-lean at commit 42df7b8}.
Its library directories and two root aggregators contain 1,172 Lean files,
436,473 physical lines, and 346,662 nonblank code lines after lexical
removal of nested block comments and line comments.
The active library alone has 429,620 physical lines in 1,153 files.
Dependencies such as mathlib are excluded. The same physical count was
verified on the VM checkout \texttt{/home/lesha/public-release-4331}.
The working \texttt{KolmogorovMathlib2-port434} checkout has the same count.
Another working branch, \texttt{KolmogorovMathlib2-merging}, contains 467,097
physical lines in 1,228 Lean files with the stated exclusions.

This article's full development is estimated at \textbf{40,000--80,000
physical lines}, or about \textbf{30,000--60,000 code lines}, with a planning
centre of 50,000--60,000 physical lines. This is an estimate, not a measured
ratio of pages to code or a guaranteed upper bound. It is approximately
9--18 percent of the public comparison snapshot. The finite certified
control core is estimated at 13,000--26,000 physical lines. Without a
proved geometric instantiation, that core gives conditional theorems.

\begin{center}
\begin{tabular}{lr}
Module family & Estimated code lines\\ \hline
Polygonal disks, placements, tilings, substitutions & 4,000--8,000\\
Local rules, crowns, positive uniform radii & 4,000--7,000\\
Compactness and infinite hierarchies & 2,000--4,000\\
Regular grouping, channels, ports, routing & 8,000--16,000\\
Decorations, origins, repaired composition & 5,000--9,000\\
STS soundness, completeness, glue dissection & 2,000--4,000\\
FLC neighbourhoods and clone reduction & 3,000--6,000\\
Integration, regression statements, final API & 2,000--6,000\\ \hline
Total & 30,000--60,000
\end{tabular}
\end{center}

The comparison project uses mature computability infrastructure from
mathlib. Here the dominant cost is geometric infrastructure rather than
the finite header algebra. Current
\href{https://leanprover-community.github.io/mathlib4_docs/Mathlib/Geometry/Polygon/Basic.html}
{Polygon.Basic} is a boundary and vertex interface, and
\href{https://leanprover-community.github.io/mathlib4_docs/Mathlib/Combinatorics/Tiling/Tile.html}
{Tiling.Tile} explicitly concerns discrete tilings. Neither search result
supplies the polygonal-disk grouping and routing needed here.

Useful existing interfaces include finite dependent function types,
K\"onig's finite inverse-system lemma, finite-product compactness,
\texttt{Multiset.count}, and graph walks. Fingerprints must preserve
multiplicity; a \texttt{Finset} of neighbour types is insufficient.
Exact APIs must be checked against the pinned mathlib commit.

A suitable first milestone is the finite control system, its repaired
composition theorem, and explicit regression statements for the old
failure. The next milestone establishes the regular geometric certificate.
Only after it is instantiated should the universal final theorem names
be marked proved. The detailed dependency graph is part of the HTML
blueprint; green kernel nodes currently cover only the header algebra.

\end{document}
